\chapter{级数与卷积}
\label{chap:series_conv}
\markboth{级数与卷积}{级数与卷积}

\begin{remark}[本章定位]\label{rem:series_chapter_positioning}
本章建立\textbf{级数}与\textbf{卷积}的完整理论。级数是函数在可列个基本单元上的分解：泰勒级数（局域导数谱）、洛朗级数（含负幂次的复展开）、傅里叶级数（谐波格上的全局频率分解）与周期级数（以振幅、频率、相位、周期为分量的复变级数，周期作一级指标）。卷积是传入两个函数与时延、计算时延乘积的函数算子，其频域形式——\textbf{相位卷积}、\textbf{振幅卷积}、\textbf{频率卷积}与\textbf{周期卷积}——将“标量与矢量”章的相似度/趋势谱系提升为频率对之间的复数矩阵。

级数是相似度与趋势在频域层面的最终归宿：傅里叶系数与玻尔系数由积分（趋势）给出，泰勒系数由导数（方向）给出。一切经典积分变换——傅里叶变换、梅林变换、拉普拉斯变换、小波变换、Chirp-Z 变换——都是周期变换在特定频率字典或核下的特殊情形。

本章最后建立\textbf{周期卷积}的正式定义：周期卷积以振幅、频率、相位、周期四个参数为算子，保留全部频率对与全部周期对的交互，是卷积理论五级阶梯（卷积 $\to$ 相位卷积/振幅卷积/频率卷积 $\to$ 趋势卷积 $\to$ 周期级数 $\to$ 周期卷积）的顶点。

\textbf{本章七主题导引}：平移——泰勒级数的局域平移展开与傅里叶级数的全局平移周期；缩放——振幅（缩放）与周期（定长）；旋转——相位（旋转）与频率（定序）；守恒——帕塞瓦尔恒等式（能量守恒）与玻尔广义帕塞瓦尔恒等式；对称——洛朗级数的正负幂对称与傅里叶基的正交对称；闭合——周期级数系数的唯一性（玻尔定理）；周期——周期级数的四元组结构与周期卷积。
\end{remark}

\section{\texorpdfstring{级数的基本概念}{级数的基本概念}}

\begin{definition}[级数与收敛]\label{def:series_basic}
设 $\{a_k\}_{k=0}^{\infty}$ 为数列。称形式和 $\sum_{k=0}^{\infty} a_k$ 为\textbf{级数}，其部分和为 $S_n = \sum_{k=0}^{n} a_k$。若 $\lim_{n\to\infty} S_n = S$ 存在，则级数\textbf{收敛}于 $S$，否则\textbf{发散}。若 $\sum_k |a_k|$ 收敛，则称原级数\textbf{绝对收敛}；收敛但不绝对收敛者称为\textbf{条件收敛}。
\end{definition}

\begin{definition}[函数项级数与一致收敛]\label{def:series_of_functions}
设 $\{u_k(t)\}$ 为 $I \subseteq \mathbb{R}$ 上的函数列。函数项级数 $\sum_k u_k(t)$ 称为在 $I$ 上\textbf{点态收敛}于 $f$，若对每个 $t \in I$ 有 $f(t) = \sum_k u_k(t)$；称为\textbf{一致收敛}于 $f$，若
\[
\lim_{n \to \infty} \sup_{t \in I} \Big| f(t) - \sum_{k=0}^{n} u_k(t) \Big| = 0.
\]
\end{definition}

\begin{proposition}[比值判别法与根值判别法]\label{prop:ratio_root_tests}
设 $a_k \neq 0$。若下列极限之一存在：
\[
\lim_{k\to\infty} \left| \frac{a_{k+1}}{a_k} \right| = \rho,
\qquad\text{或}\qquad
\lim_{k\to\infty} \sqrt[k]{|a_k|} = \rho,
\]
则
\[
\begin{aligned}
\rho < 1 &\implies \sum_k a_k \ \text{绝对收敛}, \\
\rho > 1 &\implies \sum_k a_k \ \text{发散}, \\
\rho = 1 &\implies \text{判别失效}.
\end{aligned}
\]
\end{proposition}

\begin{proof}
\textbf{比值判别法}：设 $\rho < 1$。取 $r$ 满足 $\rho < r < 1$。存在 $N$ 使 $k \ge N$ 时 $|a_{k+1}/a_k| < r$。归纳得 $|a_k| \le |a_N| r^{k-N}$。由几何级数收敛性，$\sum_k |a_k|$ 收敛。若 $\rho > 1$，取 $1 < r < \rho$，存在 $N$ 使 $k \ge N$ 时 $|a_{k+1}/a_k| > r$，故 $|a_k| \to \infty$，级数发散。$\rho = 1$ 时失效的反例：$\sum 1/k$（发散）与 $\sum 1/k^2$（收敛）的比值极限均为 $1$。

\textbf{根值判别法}：类似，利用 $\sqrt[k]{|a_k|} < r$ 或 $> r$ 的情形。
\end{proof}

\begin{theorem}[魏尔斯特拉斯 M-判别法]\label{thm:weierstrass_mtest}
设 $\sum_k M_k$ 为收敛的正项级数，且 $|u_k(t)| \le M_k$ 对所有 $t \in I$ 与 $k$ 成立。则 $\sum_k u_k(t)$ 在 $I$ 上一致且绝对收敛。
\end{theorem}

\begin{proof}
由 $\sum M_k$ 收敛，对任意 $\varepsilon > 0$ 存在 $N$ 使 $\sum_{k>n} M_k < \varepsilon$（$n \ge N$）。对任意 $t \in I$ 和 $m > n \ge N$，
\[
\Big| \sum_{k=n+1}^{m} u_k(t) \Big| \le \sum_{k=n+1}^{m} |u_k(t)| \le \sum_{k=n+1}^{m} M_k < \varepsilon.
\]
故部分和列在 $I$ 上一致柯西，从而一致收敛；绝对收敛由 $\sum_k |u_k(t)| \le \sum_k M_k < \infty$ 给出。
\end{proof}

\begin{theorem}[一致收敛的逐项积分]\label{thm:uniform_termwise_integration}
设 $\sum_k u_k$ 在 $[a,b]$ 上一致收敛于 $f$，且每个 $u_k$ 在 $[a,b]$ 上可积。则
\[
\int_a^b f(t)\,dt = \sum_{k=0}^{\infty} \int_a^b u_k(t)\,dt.
\]
\end{theorem}

\begin{proof}
记 $S_n = \sum_{k=0}^{n} u_k$。由一致收敛，
\[
\Big| \int_a^b f - \int_a^b S_n \Big| \le \int_a^b |f - S_n| \le (b-a)\, \sup_{t \in [a,b]} |f(t) - S_n(t)| \to 0.
\]
又 $\int_a^b S_n = \sum_{k=0}^{n} \int_a^b u_k$（有限和的积分等于积分的有限和），取极限即得。
\end{proof}

\begin{theorem}[一致收敛的逐项求导]\label{thm:uniform_termwise_differentiation}
设每个 $u_k$ 在 $I$ 上连续可导，$\sum_k u_k$ 在 $I$ 上点态收敛于 $f$，且 $\sum_k u_k'$ 在 $I$ 上一致收敛。则 $f$ 可导且
\[
f'(t) = \sum_{k=0}^{\infty} u_k'(t).
\]
\end{theorem}

\begin{proof}
设 $\sum_k u_k' = g$（一致收敛）。对任意 $a, t \in I$，由定理 \ref{thm:uniform_termwise_integration}，
\[
\int_a^t g(s)\,ds = \sum_{k=0}^{\infty} \int_a^t u_k'(s)\,ds = \sum_{k=0}^{\infty} \big( u_k(t) - u_k(a) \big) = f(t) - f(a).
\]
由微积分基本定理，$f'(t) = g(t) = \sum_k u_k'(t)$。
\end{proof}

\begin{theorem}[柯西-阿达马收敛半径公式]\label{thm:cauchy_hadamard}
幂级数 $\sum_k a_k (z - z_0)^k$ 的收敛半径为
\[
R = \frac{1}{\limsup_{k\to\infty} \sqrt[k]{|a_k|}},
\]
（约定 $1/0 = \infty$，$1/\infty = 0$）：$|z - z_0| < R$ 时绝对收敛，$|z - z_0| > R$ 时发散。
\end{theorem}

\begin{proof}
设 $L = \limsup_k \sqrt[k]{|a_k|}$。由根值判别法，
\[
\limsup_{k} \sqrt[k]{|a_k|\,|z-z_0|^k} = |z-z_0| \cdot L.
\]
当 $|z-z_0| L < 1$ 即 $|z-z_0| < 1/L = R$ 时，根值小于 $1$，级数绝对收敛。当 $|z-z_0| L > 1$ 时，根值大于 $1$，通项不趋于零，级数发散。
\end{proof}

\begin{property}[收敛级数的基本性质]\label{prop:series_properties}
\begin{enumerate}[label=(\roman*)]
    \item \textbf{线性性}：若 $\sum_k a_k = A$，$\sum_k b_k = B$，则 $\sum_k (a_k + b_k) = A + B$；
    \item \textbf{通项趋于零}：若 $\sum_k a_k$ 收敛，则 $a_k \to 0$；
    \item \textbf{绝对收敛蕴含收敛}：$\sum_k |a_k|$ 收敛 $\implies$ $\sum_k a_k$ 收敛。
\end{enumerate}
\end{property}

\begin{proof}
(i) 由部分和极限的线性性。(ii) $a_n = S_n - S_{n-1} \to S - S = 0$。(iii) 部分和列 $\{\sum_{k=0}^{n} a_k\}$ 是柯西列：对 $m > n$，$|\sum_{k=n+1}^m a_k| \le \sum_{k=n+1}^m |a_k|$，而 $\sum |a_k|$ 收敛保证右端可任意小。
\end{proof}

\begin{example}[几何级数：级数的原型]\label{ex:geometric_series}
\[
\sum_{k=0}^{\infty} r^k = \frac{1}{1-r}, \qquad |r| < 1.
\]
几何级数是幂级数展开的基本构件：泰勒展开与洛朗展开的正负幂部分都经由几何级数导出。
\end{example}

\begin{proposition}[柯西收敛准则]\label{prop:cauchy_criterion_series}
实数列级数 $\sum_{k=0}^{\infty} a_k$ 收敛的充要条件是：对任意 $\varepsilon > 0$，存在 $N \in \mathbb{N}$，使得对所有 $m > n \ge N$，
\[
\Big| \sum_{k=n+1}^{m} a_k \Big| < \varepsilon.
\]
\end{proposition}

\begin{proof}
（$\Rightarrow$）设 $\sum_k a_k = S$，部分和 $S_n \to S$。对 $\varepsilon > 0$ 取 $N$ 使 $n \ge N$ 时 $|S_n - S| < \varepsilon/2$。则对 $m > n \ge N$，
\[
\Big| \sum_{k=n+1}^{m} a_k \Big| = |S_m - S_n| \le |S_m - S| + |S - S_n| < \varepsilon.
\]
（$\Leftarrow$）条件表明部分和列 $\{S_n\}$ 是柯西列。由实数完备性（“数系”章定理 \ref{thm:real_completeness}），$\{S_n\}$ 收敛，即 $\sum_k a_k$ 收敛。
\end{proof}

\begin{proposition}[比较判别法]\label{prop:comparison_test}
设 $\sum_k b_k$ 为收敛的正项级数。若存在 $N$ 使 $k \ge N$ 时 $|a_k| \le b_k$，则 $\sum_k a_k$ 绝对收敛。反之，若 $\sum_k b_k$ 发散且 $k \ge N$ 时 $0 \le b_k \le a_k$，则 $\sum_k a_k$ 发散。
\end{proposition}

\begin{proof}
设 $\sum_k b_k$ 收敛。对 $m > n \ge N$，
\[
\sum_{k=n+1}^{m} |a_k| \le \sum_{k=n+1}^{m} b_k.
\]
由 $\sum_k b_k$ 的柯西准则（命题 \ref{prop:cauchy_criterion_series}），右端可任意小，故 $\sum_k |a_k|$ 收敛。第二部分由比较不等式反向即得。
\end{proof}

\begin{proposition}[积分判别法]\label{prop:integral_test}
设 $f: [1, \infty) \to [0, \infty)$ 单调递减且连续。则级数 $\sum_{k=1}^{\infty} f(k)$ 与广义积分 $\int_1^{\infty} f(x)\,dx$ 同敛散。
\end{proposition}

\begin{proof}
对 $k \le x \le k+1$，由 $f$ 单调递减得 $f(k+1) \le f(x) \le f(k)$。积分得
\[
f(k+1) \le \int_k^{k+1} f(x)\,dx \le f(k).
\]
对 $k = 1, \dots, n-1$ 求和：
\[
\sum_{k=2}^{n} f(k) \le \int_1^n f(x)\,dx \le \sum_{k=1}^{n-1} f(k).
\]
故部分和与积分同步有界或同步无界，级数与积分同敛散。
\end{proof}

\begin{property}[绝对收敛蕴含无条件收敛]\label{prop:absolute_unconditional}
若 $\sum_k |a_k|$ 收敛，则对任意重排 $\sigma: \mathbb{N} \to \mathbb{N}$（双射），$\sum_k a_{\sigma(k)}$ 收敛且等于 $\sum_k a_k$。
\end{property}

\begin{proof}
设 $S = \sum_k a_k$。由绝对收敛，给定 $\varepsilon > 0$ 存在 $N_1$ 使 $\sum_{k > N_1} |a_k| < \varepsilon/2$；又存在 $N_2$ 使 $|\sum_{k=0}^{N_2} a_k - S| < \varepsilon/2$。取 $N = \max(N_1, N_2)$。设 $M$ 充分大，使 $\{0, 1, \dots, N\} \subseteq \{\sigma(0), \dots, \sigma(M)\}$。对 $M' > M$，
\[
\Big| \sum_{k=0}^{M'} a_{\sigma(k)} - S \Big| \le \Big| \sum_{k=0}^{N} a_k - S \Big| + \sum_{\substack{k \le M' \\ \sigma(k) > N}} |a_{\sigma(k)}| < \frac{\varepsilon}{2} + \frac{\varepsilon}{2} = \varepsilon.
\]
故重排后的部分和收敛到 $S$。
\end{proof}

\begin{example}[条件收敛的重排可以改变和]\label{ex:riemann_rearrangement}
交错调和级数 $\sum_{k=1}^{\infty} \frac{(-1)^{k+1}}{k} = \ln 2$ 条件收敛。黎曼重排定理断言：对任意 $c \in \mathbb{R} \cup \{\pm\infty\}$，存在重排 $\sigma$ 使 $\sum_k a_{\sigma(k)}$ 收敛到 $c$。这一现象与绝对收敛（性质 \ref{prop:absolute_unconditional}）形成鲜明对比。
\end{example}

\begin{figure}[htbp]
\centering
\begin{tikzpicture}[scale=0.9,>=stealth]
  \draw[->] (-0.2,0) -- (6.5,0) node[right, font=\small] {$k$};
  \draw[->] (0,-0.2) -- (0,3.2) node[above, font=\small] {$a_k$};

  \foreach \k/\h in {0.5/2.6, 1.1/1.8, 1.7/1.3, 2.3/0.95, 2.9/0.72, 3.5/0.58, 4.1/0.48, 4.7/0.40, 5.3/0.34, 5.9/0.30} {
    \fill[blue!40] (\k-0.18,0) rectangle (\k+0.18,\h);
    \draw[blue!70!black] (\k-0.18,0) rectangle (\k+0.18,\h);
  }
  \node[font=\footnotesize, blue!70!black] at (3.2,2.6) {$\sum_k a_k$ 收敛};

  \begin{scope}[yshift=-3.8cm]
    \draw[->] (-0.2,0) -- (6.5,0) node[right, font=\small] {$k$};
    \draw[->] (0,-0.2) -- (0,2.6) node[above, font=\small] {$a_k$};
    \foreach \k/\h in {0.5/2.2, 1.1/1.6, 1.7/1.15, 2.3/0.85, 2.9/0.65, 3.5/0.52, 4.1/0.43, 4.7/0.36, 5.3/0.31, 5.9/0.27} {
      \fill[red!30] (\k-0.18,0) rectangle (\k+0.18,\h);
      \draw[red!70!black] (\k-0.18,0) rectangle (\k+0.18,\h);
    }
    \node[font=\footnotesize, red!70!black] at (3.2,2.2) {$\sum_k a_k$ 发散};
  \end{scope}
\end{tikzpicture}
\caption{级数收敛（上）与发散（下）的直方图对比。收敛时，柱高随 $k$ 增大而快速下降至零；发散时，柱高下降过慢，部分和持续增长。}
\label{fig:series_convergence_intuition}
\end{figure}

\begin{remark}[级数收敛的“速度层次”]\label{rem:convergence_speed_hierarchy}
级数的收敛速度分为三个层次：
\begin{enumerate}[label=(\roman*)]
  \item \textbf{几何级数速率}：$a_k = O(r^k)$（$0 < r < 1$），如 $\sum_k r^k$、$\sum_k k^{-k}$。
  \item \textbf{幂律速率}：$a_k = O(k^{-s})$（$s > 1$），如 $\sum_k k^{-2}$、$\sum_k k^{-s}$。这是调和分析中傅里叶系数衰减的典型速率（“函数与极限”章泰勒级数光滑性）。
  \item \textbf{对数速率}：$a_k = O(1/(k \log^{1+\varepsilon} k))$，如 $\sum_k 1/(k \log^2 k)$。
\end{enumerate}
收敛速率与函数光滑性、傅里叶系数衰减、泰勒余项估计直接相关（定理 \ref{thm:taylor_lagrange}、\cref{prop:amplitude_series_properties}(v)）。
\end{remark}

\section{\texorpdfstring{泰勒级数：局域导数展开}{泰勒级数}}

\begin{definition}[泰勒级数]\label{def:taylor_series}
设 $f$ 在 $x_0$ 的邻域内无穷次可导。$f$ 在 $x_0$ 处的\textbf{泰勒级数}为
\[
f(x) \sim \sum_{k=0}^{\infty} \frac{f^{(k)}(x_0)}{k!}(x - x_0)^k.
\]
其前 $n+1$ 项部分和称为 $n$ 阶\textbf{泰勒多项式} $P_n(x)$。若该级数在 $x_0$ 的某邻域内收敛于 $f$，则称 $f$ 在该邻域内\textbf{实解析}。
\end{definition}

\begin{theorem}[带拉格朗日余项的泰勒公式]\label{thm:taylor_lagrange}
设 $f$ 在 $[x_0, x]$（或 $[x, x_0]$）上 $n+1$ 次连续可导，且在开区间内 $n+1$ 次可导。则存在 $\xi$ 介于 $x_0$ 与 $x$ 之间，使
\[
f(x) = \sum_{k=0}^{n} \frac{f^{(k)}(x_0)}{k!}(x-x_0)^k + \frac{f^{(n+1)}(\xi)}{(n+1)!}(x-x_0)^{n+1}.
\]
\end{theorem}

\begin{proof}
记余项 $R_n(x) = f(x) - P_n(x)$。定义辅助函数
\[
F(t) = f(x) - \sum_{k=0}^{n} \frac{f^{(k)}(t)}{k!}(x-t)^k - R_n(x)\cdot\frac{(x-t)^{n+1}}{(x-x_0)^{n+1}}.
\]
代入得 $F(x_0) = f(x) - P_n(x) - R_n(x) = 0$，且 $F(x) = 0$。由罗尔定理，存在 $\xi$ 使 $F'(\xi) = 0$。逐项求导时中间项两两抵消，仅余首末两项：
\[
F'(t) = -\frac{f^{(n+1)}(t)}{n!}(x-t)^n + \frac{(n+1) R_n(x)}{(x-x_0)^{n+1}}(x-t)^n.
\]
令 $t = \xi$ 即得余项公式。
\end{proof}

\begin{corollary}[泰勒级数收敛于函数的充分条件]\label{cor:taylor_analytic}
若存在常数 $M > 0$ 使得 $f$ 在区间 $I$ 上各阶导数满足 $|f^{(k)}(x)| \le M$（$x \in I$，$k \ge 0$），则 $f$ 在 $I$ 内每一点的泰勒级数收敛于 $f$。
\end{corollary}

\begin{proof}
拉格朗日余项满足
\[
|R_n(x)| \le \frac{M}{(n+1)!}|x - x_0|^{n+1} \to 0 \qquad (n \to \infty).
\]
\end{proof}

\begin{example}[基本初等函数的泰勒级数]\label{ex:taylor_examples}
\[
e^x = \sum_{k=0}^{\infty} \frac{x^k}{k!}, \qquad
\sin x = \sum_{k=0}^{\infty} \frac{(-1)^k x^{2k+1}}{(2k+1)!}, \qquad
\cos x = \sum_{k=0}^{\infty} \frac{(-1)^k x^{2k}}{(2k)!},
\]
\[
\ln(1+x) = \sum_{k=1}^{\infty} \frac{(-1)^{k+1} x^k}{k} \quad (|x| < 1), \qquad
(1+x)^{\alpha} = \sum_{k=0}^{\infty} \binom{\alpha}{k} x^k \quad (|x| < 1).
\]
\end{example}

\begin{proposition}[泰勒展开的系数唯一性]\label{prop:taylor_coefficient_uniqueness}
设 $f$ 在 $x_0$ 的邻域内可展开为幂级数 $f(x) = \sum_{k=0}^{\infty} a_k (x-x_0)^k$。则 $a_k = f^{(k)}(x_0)/k!$。
\end{proposition}

\begin{proof}
对级数逐项求导 $k$ 次并在 $x = x_0$ 处取值。由幂级数在收敛圆内可逐项求导（定理 \ref{thm:uniform_termwise_differentiation}），得 $f^{(k)}(x_0) = k!\,a_k$，即 $a_k = f^{(k)}(x_0)/k!$。
\end{proof}

\begin{corollary}[泰勒多项式的逼近阶]\label{cor:taylor_approximation_order}
设 $f \in C^{n+1}(I)$，$x_0 \in I$。则泰勒多项式 $P_n(x)$ 满足
\[
f(x) - P_n(x) = O(|x - x_0|^{n+1}) \qquad (x \to x_0).
\]
\end{corollary}

\begin{proof}
由拉格朗日余项公式（定理 \ref{thm:taylor_lagrange}），$|R_n(x)| \le \frac{M_{n+1}}{(n+1)!}|x-x_0|^{n+1}$，其中 $M_{n+1} = \sup_{|t-x_0| \le |x-x_0|}|f^{(n+1)}(t)| < \infty$（因 $f^{(n+1)}$ 连续）。
\end{proof}

\begin{figure}[htbp]
\centering
\begin{tikzpicture}[scale=1.0,>=stealth]
  \draw[->] (-3.2,0) -- (3.2,0) node[right, font=\small] {$x$};
  \draw[->] (0,-0.5) -- (0,2.6) node[above, font=\small] {$y$};

  \draw[thick, blue!70!black, domain=-2.8:2.8, samples=120]
    plot (\x, {1.0 + 0.5*\x + 0.15*\x*\x - 0.02*\x*\x*\x});

  \draw[thick, red, dashed, domain=-2.0:2.0, samples=80]
    plot (\x, {1.0 + 0.5*\x});
  \node[red, font=\footnotesize] at (-2.4,0.1) {$P_1$};

  \draw[thick, green!60!black, dashed, domain=-2.4:2.4, samples=80]
    plot (\x, {1.0 + 0.5*\x + 0.15*\x*\x});
  \node[green!60!black, font=\footnotesize] at (2.4,2.2) {$P_2$};

  \draw[thick, orange!80!black, dashed, domain=-2.8:2.8, samples=100]
    plot (\x, {1.0 + 0.5*\x + 0.15*\x*\x - 0.02*\x*\x*\x});
  \node[orange!80!black, font=\footnotesize] at (-2.6,0.5) {$P_3$};

  \filldraw[black] (0,1.0) circle (2pt);
  \node[font=\footnotesize, below left] at (0,1.0) {$x_0$};

  \node[font=\small\bfseries] at (0,-0.9) {随着阶数 $n$ 增大，$P_n$ 在 $x_0$ 附近逼近 $f$};
\end{tikzpicture}
\caption{泰勒多项式的逐阶逼近。$P_1$（线性）在 $x_0$ 附近逼近 $f$；$P_2$（二次）加入曲率修正；$P_3$（三次）进一步修正。阶数越高，逼近区间越宽。}
\label{fig:taylor_approximation}
\end{figure}

\begin{remark}[泰勒展开与傅里叶展开的对偶]\label{rem:taylor_fourier_duality}
泰勒展开与傅里叶展开构成一对\textbf{局域-全局}对偶：
\begin{center}
\begin{tabular}{@{}p{3.2cm}p{4.2cm}p{4.2cm}@{}}
\toprule
 & 泰勒展开 & 傅里叶展开 \\
\midrule
展开中心 & 单点 $x_0$ & 周期区间 $[0,T]$ \\
系数来源 & 导数 $f^{(k)}(x_0)$ & 积分 $c_n = \frac{1}{T}\int_0^T f(t) e^{-i\omega_n t} dt$ \\
信息类型 & 方向（局域导数） & 趋势（全局积分） \\
收敛条件 & 实解析 & $L^2$ 可积 + 狄利克雷条件 \\
基函数 & $(x-x_0)^k$（多项式） & $e^{i\omega_n t}$（复指数） \\
\bottomrule
\end{tabular}
\end{center}
这一对偶是“标量与矢量”章方向-趋势对偶在函数展开层面的体现（注 \ref{rem:fourier_coefficient_duality}）。
\end{remark}

\section{\texorpdfstring{洛朗级数：含负幂次的复展开}{洛朗级数}}

\begin{definition}[洛朗级数]\label{def:laurent_series}
设 $z_0 \in \mathbb{C}$。形如
\[
f(z) = \sum_{n=-\infty}^{\infty} a_n (z - z_0)^n
\]
的级数称为以 $z_0$ 为中心的\textbf{洛朗级数}。正幂部分在圆盘 $|z-z_0| < R$ 内收敛，负幂部分经换元 $w = 1/(z-z_0)$ 在 $|z-z_0| > r$ 上收敛；共同收敛域为圆环 $r < |z-z_0| < R$。
\end{definition}

\begin{lemma}[洛朗系数的唯一性]\label{lem:laurent_uniqueness}
若 $\sum_{n} a_n (z-z_0)^n = \sum_{n} b_n (z-z_0)^n$ 在某圆环上成立，则 $a_n = b_n$ 对所有 $n \in \mathbb{Z}$ 成立。
\end{lemma}

\begin{proof}
令 $d_n = a_n - b_n$。对任意 $m \in \mathbb{Z}$，将恒等式乘以 $(z-z_0)^{-m-1}$ 沿围绕 $z_0$ 的圆积分。由
\[
\oint_\gamma (z-z_0)^{n-m-1}\,dz = 2\pi i\, \delta_{nm},
\]
得 $d_m \cdot 2\pi i = 0$，故 $d_m = 0$。
\end{proof}

\begin{theorem}[洛朗展开定理]\label{thm:laurent_expansion}
设 $f$ 在圆环 $r < |z - z_0| < R$ 内解析。则 $f$ 在该圆环内可唯一展开为洛朗级数，系数为
\[
a_n = \frac{1}{2\pi i} \oint_{\gamma} \frac{f(\zeta)}{(\zeta - z_0)^{n+1}}\, d\zeta, \qquad n \in \mathbb{Z},
\]
其中 $\gamma$ 为圆环内围绕 $z_0$ 的任一正向闭曲线。
\end{theorem}

\begin{proof}
由柯西积分公式，对圆环内任一点 $z$ 与两同心圆 $\gamma_1, \gamma_2$，
\[
f(z) = \frac{1}{2\pi i}\oint_{\gamma_1} \frac{f(\zeta)}{\zeta - z}\,d\zeta - \frac{1}{2\pi i}\oint_{\gamma_2} \frac{f(\zeta)}{\zeta - z}\,d\zeta.
\]
在 $\gamma_1$ 上用几何级数展开得正幂部分；在 $\gamma_2$ 上类似展开得负幂部分。唯一性由引理 \ref{lem:laurent_uniqueness} 给出。
\end{proof}

\begin{proposition}[旋转周期函数的洛朗展开]\label{prop:laurent_rotation_period}
若 $f$ 满足旋转周期 $f(e^{i\theta} z) = f(z)$（具有周期轨道的群 $C_m$，即 $\theta = 2\pi/m$），则其以 $0$ 为中心的洛朗级数仅含 $n \equiv 0 \pmod{m}$ 的项。
\end{proposition}

\begin{proof}
代入 $z \mapsto e^{i\theta} z$，由系数唯一性，$a_n (e^{in\theta} - 1) = 0$。当 $e^{i\theta}$ 为 $m$ 次本原单位根时，$e^{in\theta} = 1 \iff m \mid n$。
\end{proof}

\begin{example}[$e^{1/z}$ 的洛朗展开]\label{ex:laurent_example}
\[
e^{1/z} = \sum_{n=0}^{\infty} \frac{z^{-n}}{n!}, \qquad |z| > 0.
\]
\end{example}

\begin{proposition}[洛朗级数的奇点分类与主部]\label{prop:laurent_singularity}
设 $f$ 在 $0 < |z-z_0| < R$ 上解析，洛朗展开为 $f(z) = \sum_{n=-\infty}^{\infty} a_n (z-z_0)^n$。称
\[
\text{主部} = \sum_{n=1}^{\infty} a_{-n} (z-z_0)^{-n}
\]
为洛朗展开的主部，则：
\begin{enumerate}[label=(\roman*)]
  \item 可去奇点 $\iff$ 主部为零；
  \item $m$ 阶极点 $\iff$ 主部为非零的有限和，且最高负幂次为 $m$；
  \item 本质奇点 $\iff$ 主部有无穷多项。
\end{enumerate}
\end{proposition}

\begin{proof}
逐条由洛朗展开的负幂部分定义直接得出。
\end{proof}

\begin{example}[三类奇点的洛朗展开对比]\label{ex:three_singularity_types}
\begin{enumerate}[label=(\roman*)]
  \item \textbf{可去奇点}：$f(z) = \frac{\sin z}{z}$，在 $z_0 = 0$ 处洛朗展开为 $1 - z^2/6 + \cdots$，主部为零。
  \item \textbf{极点}：$f(z) = \frac{1}{z^2(z-1)}$，在 $z_0 = 0$ 处主部为 $-1 - z - z^2$，故为 $2$ 阶极点。
  \item \textbf{本质奇点}：$f(z) = e^{1/z}$，在 $z_0 = 0$ 处主部为 $\sum_{n=1}^{\infty} z^{-n}/n!$，有无穷多项。
\end{enumerate}
\end{example}

\begin{theorem}[留数定理与傅里叶系数的对偶]\label{thm:residue_fourier_duality}
设 $f$ 在圆环 $r < |z - z_0| < R$ 上解析，洛朗展开为 $f(z) = \sum_{n \in \mathbb{Z}} a_n (z-z_0)^n$。则 $a_n$ 由留数表示为
\[
a_n = \operatorname{Res}\!\left(\frac{f(z)}{(z-z_0)^{n+1}}, z_0\right).
\]
特别地，$a_{-1} = \operatorname{Res}(f, z_0)$ 恰为洛朗展开中的 $(z-z_0)^{-1}$ 系数。
\end{theorem}

\begin{proof}
由洛朗系数公式（定理 \ref{thm:laurent_expansion}）与留数定义（“微分方程与解析”章定义 \ref{def:residue}）直接对比即得。
\end{proof}

\begin{figure}[htbp]
\centering
\begin{tikzpicture}[scale=1.1,>=stealth]
  \draw[gray!60, thin] (0,0) circle (2.2);
  \draw[gray!60, thin] (0,0) circle (1.2);
  \filldraw[red!70!black] (0,0) circle (2.5pt);
  \node[red!70!black, font=\footnotesize, below left=1pt] at (0,0) {$z_0$};

  \fill[blue!10, opacity=0.5] (0,0) circle (1.2);
  \fill[green!10, opacity=0.4] (0,0) circle (2.2);
  \fill[white, opacity=1.0] (0,0) circle (1.2);
  \fill[red!5, opacity=0.7] (0,0) circle (1.2);

  \node[font=\scriptsize, blue!60!black] at (0,0.6) {收敛};
  \node[font=\scriptsize, blue!60!black] at (0,0.3) {圆盘};
  \node[font=\scriptsize, green!60!black] at (0,1.8) {圆环 $r < |z - z_0| < R$};

  \draw[->, thick, blue] (1.3,0.5) -- (1.8,0.5) node[right, font=\footnotesize, black] {正幂收敛域};

  \draw[gray!70, dashed] (0,0) -- (0,1.2);
  \node[gray!70, font=\footnotesize] at (0.35,0.6) {$r$};
  \draw[gray!70, dashed] (0,0) -- ({cos(60)*2.2}, {sin(60)*2.2});
  \node[gray!70, font=\footnotesize] at ({cos(60)*1.1}, {sin(60)*1.1 + 0.15}) {$R$};
\end{tikzpicture}
\caption{洛朗级数的收敛圆环。正幂部分在以 $z_0$ 为中心的圆盘内收敛，负幂部分在无穷远点附近的区域收敛；二者共同收敛域为圆环 $r < |z - z_0| < R$。}
\label{fig:laurent_annulus}
\end{figure}

\section{\texorpdfstring{傅里叶级数：谐波格上的全局分解}{傅里叶级数}}

\begin{definition}[傅里叶级数（复形式与实形式）]\label{def:fourier_series_foundations}
设 $f$ 以 $T$ 为周期且在 $[0,T]$ 上可积。$f$ 的\textbf{傅里叶级数}（复形式）为
\[
f(t) \sim \sum_{n=-\infty}^{\infty} c_n e^{i 2\pi n t / T}, \qquad
c_n = \frac{1}{T}\int_0^T f(t) e^{-i 2\pi n t / T}\,dt.
\]
实形式为
\[
f(t) \sim \frac{a_0}{2} + \sum_{n=1}^{\infty} \Big( a_n \cos\frac{2\pi n t}{T} + b_n \sin\frac{2\pi n t}{T} \Big),
\]
其中 $a_n = c_n + c_{-n}$，$b_n = i(c_n - c_{-n})$。
\end{definition}

\begin{proposition}[傅里叶基的正交性]\label{prop:fourier_orthogonality}
函数系 $\{e^{i 2\pi n t / T}\}_{n \in \mathbb{Z}}$ 在 $L^2([0,T])$ 中正交：
\[
\frac{1}{T}\int_0^T e^{i 2\pi (n-m) t / T}\,dt = \delta_{nm}.
\]
\end{proposition}

\begin{proof}
$n = m$ 时被积函数为 $1$；$n \neq m$ 时积分化为 $\frac{T}{2\pi i (n-m)}(e^{i 2\pi (n-m)} - 1) = 0$。
\end{proof}

\begin{theorem}[帕塞瓦尔恒等式（傅里叶级数）]\label{thm:parseval_fourier_classical}
设 $f \in L^2([0,T])$，傅里叶系数为 $c_n$。则
\[
\frac{1}{T}\int_0^T |f(t)|^2\,dt = \sum_{n=-\infty}^{\infty} |c_n|^2.
\]
\end{theorem}

\begin{proof}
对三角多项式由正交性（命题 \ref{prop:fourier_orthogonality}）直接成立；由傅里叶基在 $L^2$ 中的完备性经逼近推广到一般 $L^2$ 函数。
\end{proof}

\begin{example}[方波的傅里叶级数]\label{ex:fourier_square_wave}
设 $f$ 以 $T$ 为周期，在 $[0, T/2)$ 上取值 $1$，在 $[T/2, T)$ 上取值 $-1$。其傅里叶级数为
\[
f(t) \sim \frac{4}{\pi} \sum_{k=0}^{\infty} \frac{1}{2k+1} \sin\frac{2\pi (2k+1) t}{T}.
\]
\end{example}

\begin{theorem}[狄利克雷收敛定理]\label{thm:dirichlet_convergence}
设 $f$ 以 $T$ 为周期，在 $[0,T]$ 上分段 $C^1$。则 $f$ 的傅里叶级数在每一点 $t$ 收敛于
\[
\frac{f(t^+) + f(t^-)}{2},
\]
其中 $f(t^+)$、$f(t^-)$ 分别为 $f$ 在 $t$ 处的右极限与左极限。特别地，在 $f$ 的连续点处，傅里叶级数收敛于 $f(t)$。
\end{theorem}

\begin{proof}[证明思路]
将 $f$ 的傅里叶部分和写成与狄利克雷核 $D_N$ 的卷积：
\[
S_N f(t) = \frac{1}{T}\int_0^T f(t - \tau) D_N(\tau)\,d\tau.
\]
利用狄利克雷核的性质：$\frac{1}{T}\int_0^T D_N(\tau)\,d\tau = 1$，且 $D_N$ 在远离 $0$ 处正负振荡相消（黎曼–勒贝格引理）。分段 $C^1$ 条件保证 $f$ 在 $t$ 附近的左右局部行为可由一阶泰勒展开控制，从而 $S_N f(t) \to [f(t^+) + f(t^-)]/2$。严格证明见 Zygmund, \textit{Trigonometric Series}, 3rd ed., Cambridge Univ. Press, 2002, \S{}II.4。
\end{proof}

\begin{corollary}[连续点处的点态收敛]\label{cor:fourier_continuity}
若 $f$ 在 $t$ 处连续，则 $f(t^+) = f(t^-) = f(t)$，傅里叶级数在 $t$ 处收敛于 $f(t)$。
\end{corollary}

\begin{example}[吉布斯现象]\label{ex:gibbs_phenomenon}
方波 $f(t) = \operatorname{sgn}(\sin t)$ 的傅里叶部分和 $S_N f$ 在跳跃点 $t = 0$ 附近出现约 $9\%$ 的\textbf{过冲}（overshoot），且过冲幅度不随 $N \to \infty$ 衰减。这一现象称为\textbf{吉布斯现象}，说明傅里叶级数在跳跃点附近的一致收敛不成立（点态收敛仍成立，由定理 \ref{thm:dirichlet_convergence}）。吉布斯现象在数字信号处理中与“理想低通滤波”的振铃效应等价。
\end{example}

\begin{theorem}[费耶核与一致收敛]\label{thm:fejer_uniform}
设 $f$ 以 $T$ 为周期且连续。定义费耶平均
\[
\sigma_N f(t) = \frac{1}{N+1}\sum_{k=0}^{N} S_k f(t),
\]
其中 $S_k f$ 为第 $k$ 个傅里叶部分和。则 $\sigma_N f$ 在 $\mathbb{R}$ 上一致收敛于 $f$。
\end{theorem}

\begin{proof}[证明思路]
费耶平均可写成与费耶核的卷积 $\sigma_N f(t) = \frac{1}{T}\int_0^T f(t-\tau) F_N(\tau)\,d\tau$，其中
\[
F_N(\tau) = \frac{1}{N+1}\left(\frac{\sin((N+1)\tau/2)}{\sin(\tau/2)}\right)^2.
\]
费耶核是\textbf{好核}：非负、积分为 $1$、在 $\tau \neq 0$ 处一致趋于零。故对连续周期函数 $f$，由好核的一致逼近性质得 $\sigma_N f \to f$ 一致成立。严格证明见 Zygmund, \textit{Trigonometric Series}, 3rd ed., Cambridge Univ. Press, 2002, \S{}III.1。
\end{proof}

\begin{figure}[htbp]
\centering
\begin{tikzpicture}[scale=0.9,>=stealth]
  \draw[->] (-3.5,0) -- (3.5,0) node[right, font=\small] {$t$};
  \draw[->] (0,-1.6) -- (0,1.8) node[above, font=\small] {$f$};

  \draw[thick, gray!50, domain=-3.2:3.2, samples=200]
    plot (\x, {(\x>0?1:-1)});

  \draw[thick, blue!70!black, domain=-3.2:3.2, samples=200]
    plot (\x, {0.9*sin(deg(2*\x))/(2*\x) + 0.05*sin(deg(6*\x))/(6*\x) + 0.03*sin(deg(10*\x))/(10*\x)});
  \node[blue!70!black, font=\footnotesize] at (2.5,1.6) {$S_N$};

  \draw[red, thick, dashed] (-0.15,1.18) -- (0.15,1.18);
  \node[red!70!black, font=\scriptsize, above] at (0,1.18) {过冲 $\approx 9\%$};
\end{tikzpicture}
\caption{吉布斯现象。方波（灰线）的傅里叶部分和 $S_N$（蓝线）在跳跃点 $t = 0$ 附近出现过冲（红虚线），过冲幅度约为跳跃幅度的 $9\%$，且不随 $N$ 增大而消失。}
\label{fig:gibbs}
\end{figure}

\begin{property}[傅里叶级数的积分保持]\label{prop:fourier_integration_preservation}
设 $f$ 以 $T$ 为周期，在 $[0,T]$ 上可积且均值为零（即 $c_0 = 0$）。则 $f$ 的不定积分 $F(t) = \int_0^t f(\tau)\,d\tau$ 以 $T$ 为周期，且其傅里叶系数为 $C_n = c_n / (i\omega_n)$（$n \neq 0$）。
\end{property}

\begin{proof}
由 $c_0 = 0$ 得 $F(T) = \int_0^T f = T c_0 = 0$，故 $F$ 以 $T$ 为周期。$F$ 的傅里叶系数为
\[
C_n = \frac{1}{T}\int_0^T F(t) e^{-i\omega_n t}\,dt = \frac{1}{T} \cdot \frac{1}{i\omega_n}\left[ F(t) e^{-i\omega_n t} \Big|_0^T + \int_0^T f(t) e^{-i\omega_n t}\,dt \right] = \frac{c_n}{i\omega_n}.
\]
（分部积分，边界项由周期性为零。）
\end{proof}

\begin{remark}[傅里叶级数三种收敛模式]\label{rem:fourier_convergence_modes}
傅里叶级数具有三种收敛模式，各自对应不同条件与结论：
\begin{center}
\begin{tabular}{@{}p{2.5cm}p{4.0cm}p{5.0cm}@{}}
\toprule
收敛模式 & 条件 & 结论 \\
\midrule
逐点收敛 & $f$ 分段 $C^1$ & $S_N f(t) \to [f(t^+)+f(t^-)]/2$（定理 \ref{thm:dirichlet_convergence}） \\
一致收敛 & $f$ 连续 + $\sum |c_n| < \infty$ & $S_N f \to f$ 一致（费耶平均，定理 \ref{thm:fejer_uniform}） \\
$L^2$ 收敛 & $f \in L^2[0,T]$ & $\|S_N f - f\|_2 \to 0$（帕塞瓦尔 + 完备性） \\
\bottomrule
\end{tabular}
\end{center}
三种模式互不蕴含：一致收敛 $\Rightarrow$ 逐点收敛，但反之不成立（吉布斯现象）；$L^2$ 收敛与逐点收敛互不蕴含。
\end{remark}

\begin{remark}[傅里叶系数的方向-趋势对偶]\label{rem:fourier_coefficient_duality}
傅里叶系数 $c_n$ 由积分（趋势信息）给出，泰勒系数由导数（方向信息）给出，两者形成方向-趋势对偶。傅里叶级数的点态收敛与一致收敛问题可由狄利克雷核、吉布斯现象、费耶核平滑等工具处理，本章不展开。
\end{remark}

\section{\texorpdfstring{河图复变函数：周期 10 的 $\mathbb{Z}_2 \times \mathbb{Z}_5$ 分解}{河图复变函数}}
\label{sec:hetu_function}

本节建立\textbf{河图复变函数}——以 $10$ 为周期的复变函数。其数学核心是：周期 $10$ 的傅里叶分解中自动出现 $\mathbb{Z}_{10} \cong \mathbb{Z}_2 \times \mathbb{Z}_5$ 的代数分级，其中 $\mathbb{Z}_5$ 分量对应五重循环，$\mathbb{Z}_2$ 分量对应二重交替。该分级是离散周期函数傅里叶展开的纯代数事实，不依赖任何文化诠释。

\begin{definition}[$\mathbb{Z}_{10}$ 上的 $\mathbb{Z}_5$ 投影]\label{def:hetu_mapping}
定义群同态 $\varphi: \mathbb{Z}_{10} \to \mathbb{Z}_5$ 为
\[
\varphi(n) := n \bmod 5, \qquad n \in \mathbb{Z}_{10} = \{0, 1, \dots, 9\},
\]
其中约定余数 $0$ 对应标签“土”。具体地：
\[
\begin{array}{c|ccccc}
\varphi(n) & 0 & 1 & 2 & 3 & 4 \\
\hline
\text{标签} & \text{土} & \text{水} & \text{火} & \text{木} & \text{金} \\
\text{数字} & 5, 10 & 1, 6 & 2, 7 & 3, 8 & 4, 9
\end{array}
\]
另有 $\mathbb{Z}_2$ 分量为 $n$ 的奇偶性。$\varphi$ 的核为
\[
\ker \varphi = \{0, 5\} \cong \mathbb{Z}_2.
\]
\end{definition}

\begin{definition}[河图复变函数]\label{def:hetu_function}
设 $f: \mathbb{C} \to \mathbb{C}$ 满足周期条件
\[
f(z + 10) = f(z), \qquad \forall z \in \mathbb{C}.
\]
则称 $f$ 为\textbf{河图复变函数}。其傅里叶展开为
\[
f(z) = \sum_{n \in \mathbb{Z}} c_n\, e^{2\pi i n z / 10}, \qquad c_n = \frac{1}{10}\int_0^{10} f(t)\, e^{-2\pi i n t / 10}\, dt.
\]
\end{definition}

\begin{theorem}[河图傅里叶分解——$\mathbb{Z}_2 \times \mathbb{Z}_5$ 分级]\label{thm:hetu_fourier_decomposition}
设 $f$ 为河图复变函数，其傅里叶展开在 $\mathbb{C}$ 上一致收敛。则 $f$ 可唯一分解为
\[
\boxed{\,f(z) = \sum_{k=0}^{4} \Bigl( c_k + c_{k+5}\, e^{\pi i z} \Bigr)\, e^{2\pi i k z / 10}\,,}
\]
其中 $c_k$ 为傅里叶系数，$e^{\pi i z}$ 是周期为 $2$ 的 $\mathbb{Z}_2$ 因子。该分解是 $\mathbb{Z}_{10} \cong \mathbb{Z}_2 \times \mathbb{Z}_5$ 在傅里叶展开中的显式形式。
\end{theorem}

\begin{proof}
利用 $\mathbb{Z}_{10} \cong \mathbb{Z}_2 \times \mathbb{Z}_5$，将整数指标 $n \in \mathbb{Z}$ 唯一地写成
\[
n = 5\varepsilon + k, \qquad \varepsilon \in \{0, 1\}, \quad k \in \{0, 1, 2, 3, 4\}.
\]
代入傅里叶基函数：
\[
e^{2\pi i n z / 10} = e^{2\pi i (5\varepsilon + k) z / 10} = e^{\pi i \varepsilon z} \cdot e^{2\pi i k z / 10}.
\]
将 $n$ 的求和按 $\varepsilon$ 分解：
\[
f(z) = \sum_{\varepsilon \in \{0,1\}} \sum_{k=0}^{4} c_{5\varepsilon+k}\, e^{\pi i \varepsilon z}\, e^{2\pi i k z / 10}.
\]
提取 $\varepsilon = 0$ 与 $\varepsilon = 1$ 两项并合并同类 $e^{2\pi i k z / 10}$：
\[
f(z) = \sum_{k=0}^{4} \Bigl( c_k + c_{k+5}\, e^{\pi i z} \Bigr)\, e^{2\pi i k z / 10}.
\]
唯一性由傅里叶系数的唯一性（玻尔均值）与 $n \mapsto (5\varepsilon + k)$ 双射保证。
\end{proof}

\begin{definition}[$\mathbb{Z}_5$ 分量]\label{def:hetu_five_components}
在上述分解中，称
\[
f_k(z) := \Bigl( c_k + c_{k+5}\, e^{\pi i z} \Bigr)\, e^{2\pi i k z / 10}, \qquad k = 0, 1, 2, 3, 4
\]
为 $f$ 的\textbf{第 $k$ 个 $\mathbb{Z}_5$ 分量}。按定义 \ref{def:hetu_mapping} 的标签，五个分量依次对应：土、水、火、木、金。每个分量由两部分构成：$c_k\, e^{2\pi i k z / 10}$（$\varepsilon = 0$ 分量）与 $c_{k+5}\, e^{\pi i z} e^{2\pi i k z / 10}$（$\varepsilon = 1$ 分量）。
\end{definition}

\begin{proposition}[$\mathbb{Z}_5$ 分量的周期结构]\label{prop:hetu_five_periods}
五个 $\mathbb{Z}_5$ 分量 $f_0, f_1, f_2, f_3, f_4$ 的基频周期分别为：
\[
\begin{array}{c|ccccc}
k & 0 & 1 & 2 & 3 & 4 \\
\hline
\text{标签} & \text{土} & \text{水} & \text{火} & \text{木} & \text{金} \\
\text{基频周期} & 1 & 10 & 5 & 10 & 5
\end{array}
\]
即：$f_0$ 仅含常数项；$f_1, f_3$ 为周期 $10$ 的分量；$f_2, f_4$ 为周期 $5$ 的分量。
\end{proposition}

\begin{proof}
分量 $f_k$ 的基频为 $e^{2\pi i k z / 10}$，其基频周期为 $10/\gcd(k, 10)$。逐一代入：
\begin{itemize}
    \item $k = 0$：$10/\gcd(0,10) = 1$（约定 $\gcd(0,10) = 10$），故 $f_0$ 的基频部分为常数。
    \item $k = 1$：$10/\gcd(1,10) = 10$。
    \item $k = 2$：$10/\gcd(2,10) = 5$。
    \item $k = 3$：$10/\gcd(3,10) = 10$。
    \item $k = 4$：$10/\gcd(4,10) = 5$。
\end{itemize}
$\mathbb{Z}_2$ 因子 $e^{\pi i z}$ 的周期为 $2$，不改变上述周期结构（因为 $2 \mid 10$）。
\end{proof}

\begin{proposition}[河图复变函数的半周期平移]\label{prop:hetu_yin_yang_flip}
设 $f$ 为河图复变函数，$f = \sum_{k=0}^{4} f_k$ 为其 $\mathbb{Z}_5$ 分解。则半周期平移 $z \mapsto z + 5$ 使
\[
f(z + 5) = \sum_{k=0}^{4} \Bigl( (-1)^k c_k + (-1)^k c_{k+5}\, e^{\pi i z} \Bigr)\, e^{2\pi i k z / 10},
\]
即每个 $\mathbb{Z}_5$ 分量获得因子 $(-1)^k$：偶数分量（$k = 0, 2, 4$）保持不变，奇数分量（$k = 1, 3$）整体变号。
\end{proposition}

\begin{proof}
由定理 \ref{thm:hetu_fourier_decomposition}，$f(z+5) = \sum_k (c_k + c_{k+5} e^{\pi i (z+5)}) e^{2\pi i k (z+5)/10}$。分别计算：
\[
e^{\pi i (z+5)} = e^{\pi i z} e^{5\pi i} = -e^{\pi i z},\qquad
e^{2\pi i k (z+5)/10} = e^{2\pi i k z / 10}\, e^{\pi i k} = (-1)^k\, e^{2\pi i k z / 10}.
\]
代入后即得。
\end{proof}

\begin{definition}[$\mathbb{Z}_{10}$ 相位算子]\label{def:hetu_phase_operator}
\textbf{$\mathbb{Z}_{10}$ 相位算子}定义为
\[
\mathcal{O}_n^{(10)} := e^{2\pi i n / 10}, \qquad n = 0, 1, \dots, 9.
\]
在 $\mathbb{Z}_{10} \cong \mathbb{Z}_2 \times \mathbb{Z}_5$ 分解下，相位算子可写为
\[
\mathcal{O}_{5\varepsilon + k}^{(10)} = (-1)^\varepsilon\, e^{2\pi i k / 10}, \qquad \varepsilon \in \{0, 1\}, \; k \in \{0, 1, 2, 3, 4\},
\]
即每个相位算子由 $\mathbb{Z}_2$ 符号 $(-1)^\varepsilon$ 与 $\mathbb{Z}_5$ 相位 $e^{2\pi i k / 10}$ 的乘积构成。
\end{definition}

\begin{proposition}[$\mathbb{Z}_{10}$ 相位算子群的直积分解]\label{prop:hetu_phase_group}
$\mathbb{Z}_{10}$ 相位算子群 $\{\mathcal{O}_n^{(10)}\}_{n=0}^{9}$ 同构于
\[
C_{10} \cong C_2 \times C_5,
\]
其中 $C_2 = \{+1, -1\}$ 对应 $\mathbb{Z}_2$ 分量，$C_5 = \{e^{2\pi i k/5}\}_{k=0}^{4}$ 对应 $\mathbb{Z}_5$ 分量。生成元 $\mathcal{O}_1^{(10)} = e^{\pi i / 5}$ 在分解下对应 $(1, e^{2\pi i / 10})$，其阶为 $10 = \operatorname{lcm}(2, 5)$。
\end{proposition}

\begin{proof}
映射 $\mathcal{O}_{5\varepsilon + k}^{(10)} \mapsto ((-1)^\varepsilon, e^{2\pi i k / 10})$ 是群同构：由
\[
\mathcal{O}_{5\varepsilon_1 + k_1}^{(10)} \cdot \mathcal{O}_{5\varepsilon_2 + k_2}^{(10)} = \mathcal{O}_{5(\varepsilon_1 + \varepsilon_2) + (k_1 + k_2)}^{(10)},
\]
右端对应 $((-1)^{\varepsilon_1 + \varepsilon_2}, e^{2\pi i (k_1 + k_2)/10})$，恰为两分量乘积。同构为双射由 $\mathbb{Z}_{10} \cong \mathbb{Z}_2 \times \mathbb{Z}_5$ 保证。生成元阶 $= \operatorname{lcm}(2, 5) = 10$。
\end{proof}

\begin{theorem}[河图循环矩阵的 DFT 对角化]\label{thm:hetu_circulant}
设 $\mathbf{C}$ 是 $10 \times 10$ 循环矩阵，其首行生成元为 $(h_0, h_1, \dots, h_9)$，且 $h_n$ 按定义 \ref{def:hetu_mapping} 赋予 $\mathbb{Z}_2 \times \mathbb{Z}_5$ 标号。则 $\mathbf{C}$ 被 DFT 矩阵 $\mathbf{F}_{10}$ 酉对角化：
\[
\mathbf{C} = \mathbf{F}_{10}^*\, \operatorname{diag}(\hat{h}_0, \hat{h}_1, \dots, \hat{h}_9)\, \mathbf{F}_{10}, \qquad
\hat{h}_n = \sum_{j=0}^{9} h_j\, e^{-2\pi i j n / 10}.
\]
进一步，特征值 $\hat{h}_n$ 按 $\mathbb{Z}_{10} \cong \mathbb{Z}_2 \times \mathbb{Z}_5$ 分级：对应 $\varepsilon = 0$ 的五个特征值 $(\hat{h}_0, \hat{h}_1, \hat{h}_2, \hat{h}_3, \hat{h}_4)$ 属 $C_5$ 分量，对应 $\varepsilon = 1$ 的五个特征值 $(\hat{h}_5, \hat{h}_6, \hat{h}_7, \hat{h}_8, \hat{h}_9)$ 属 $\mathbb{Z}_2$ 平移分量。
\end{theorem}

\begin{proof}
由循环矩阵的 DFT 对角化（标准结果），$\mathbf{C} = \mathbf{F}_{10}^* \operatorname{diag}(\hat{h}) \mathbf{F}_{10}$。特征值 $\hat{h}_n$ 按 $n = 5\varepsilon + k$ 分为两组：$n \in \{0,1,2,3,4\}$ 对应 $\varepsilon = 0$，$n \in \{5,6,7,8,9\}$ 对应 $\varepsilon = 1$。由 $\hat{h}_{n+5} = \sum_j h_j e^{-2\pi i j (n+5)/10} = \sum_j h_j (-1)^j e^{-2\pi i j n/10}$，即 $\hat{h}_{n+5}$ 由 $h_j$ 的奇偶加权决定，与 $\hat{h}_n$ 独立。
\end{proof}

\begin{example}[河图复变函数的具体分解]\label{ex:hetu_concrete}
取 $f(z) = 3 + 2 e^{\pi i z / 5} + e^{\pi i z} + 5 e^{3\pi i z / 5} + 4 e^{2\pi i z}$。该函数的周期为 $10$，将其写为定理 \ref{thm:hetu_fourier_decomposition} 的标准形式。

先按 $n$ 顺序排列傅里叶系数：$c_0 = 3$，$c_1 = 2$，$c_5 = 1$，$c_3 = 5$，$c_4 = 4$，其余 $c_n = 0$。按 $n = 5\varepsilon + k$ 分解：
\[
f(z) = \underbrace{3}_{n=0} + \underbrace{2 e^{\pi i z / 5}}_{n=1} + \underbrace{5 e^{3\pi i z / 5}}_{n=3} + \underbrace{4 e^{2\pi i z}}_{n=4} + \underbrace{e^{\pi i z}}_{n=5}.
\]
每项分量的 $(\varepsilon, k)$ 标号与 $\mathbb{Z}_2 \times \mathbb{Z}_5$ 归属如下表：
\[
\begin{array}{c|ccccc}
n & 0 & 1 & 3 & 4 & 5 \\
\hline
(\varepsilon, k) & (0, 0) & (0, 1) & (0, 3) & (0, 4) & (1, 0) \\
\text{标签} & \text{土} & \text{水} & \text{木} & \text{金} & \text{土} \\
\text{分量} & C_5 & C_5 & C_5 & C_5 & \mathbb{Z}_2
\end{array}
\]
合并同类 $\mathbb{Z}_5$ 分量：
\[
f_0 = 3 + e^{\pi i z}, \quad f_1 = 2 e^{\pi i z/5}, \quad f_2 = 0, \quad f_3 = 5 e^{3\pi i z/5}, \quad f_4 = 4 e^{2\pi i z}.
\]
按命题 \ref{prop:hetu_five_periods}，$f_0$ 含常数项（周期 $1$）；$f_1, f_3$ 为周期 $10$；$f_4$ 为周期 $5$。
\end{example}

\begin{remark}[河图结构的代数意义]\label{rem:hetu_algebraic_meaning}
河图复变函数的数学内容在于 $\mathbb{Z}_{10} \cong \mathbb{Z}_2 \times \mathbb{Z}_5$ 的傅里叶分级：
\begin{itemize}
    \item $\mathbb{Z}_5$ 分量对应傅里叶模 $k \in \{0,1,2,3,4\}$ 的五重循环；
    \item $\mathbb{Z}_2$ 分量对应傅里叶模 $n$ 的奇偶性，即二重交替；
    \item 半周期平移 $z \mapsto z + 5$ 使 $k$ 为奇数的分量整体变号（命题 \ref{prop:hetu_yin_yang_flip}）。
\end{itemize}
河图映射 $\varphi: \mathbb{Z}_{10} \to \mathbb{Z}_5$ 是群同态，其核 $\{0, 5\} \cong \mathbb{Z}_2$。这使分级在傅里叶展开中获得严格代数实现——不是外加的编码，而是周期 $10$ 的傅里叶分解中自动出现的代数结构。
\end{remark}

\section{\texorpdfstring{振幅级数}{振幅级数}}

本节讨论相位为零、周期为 $2\pi$ 的级数形式。它是傅里叶级数在丢弃全部相位信息后得到的振幅谱，也是帕塞瓦尔能量守恒的最直接载体。

\begin{definition}[相位为零、周期为 $2\pi$ 的振幅级数]\label{def:amplitude_series}
设 $A_n \ge 0$（$n \in \mathbb{Z}$）为振幅序列。形式级数
\[
\mathcal{A}_f(t) := \sum_{n=-\infty}^{\infty} A_n e^{i n t}
\]
称为\textbf{相位为零、周期为 $2\pi$ 的振幅级数}。其实形式为
\[
\mathcal{A}_f^{\mathbb{R}}(t) := \sum_{n=-\infty}^{\infty} A_n \cos(nt),
\]
此时要求 $A_{-n}=A_n$。若需强调生成函数形式，亦记
\[
\mathfrak{A}_f(s) := \sum_{n=-\infty}^{\infty} A_n s^n.
\]
振幅级数丢弃全部相位信息，只保留各次谐波的振幅分布。当 $\sum_n A_n < \infty$ 时，$\mathcal{A}_f$ 一致绝对收敛；一般情形下按形式级数或分布意义理解。
\end{definition}

\begin{proposition}[振幅级数的玻尔均值系数公式]\label{prop:amplitude_bohr_mean}
若振幅级数 $\mathcal{A}_f(t)=\sum_n A_n e^{int}$ 一致绝对收敛，则对每个 $n \in \mathbb{Z}$，
\[
A_n = M\{\mathcal{A}_f(t)e^{-int}\}
= \lim_{T\to\infty}\frac{1}{2T}\int_{-T}^{T}\mathcal{A}_f(t)e^{-int}\,dt.
\]
\end{proposition}

\begin{proof}
由一致绝对收敛性，玻尔均值可与求和交换：
\[
M\{\mathcal{A}_f(t)e^{-int}\}
= \sum_{m\in\mathbb{Z}} A_m M\{e^{i(m-n)t}\}.
\]
由引理 \ref{lem:amplitude_exponential_mean}，
\[
M\{e^{i(m-n)t}\} = \delta_{mn},
\]
故仅 $m=n$ 项保留，得 $A_n$。
\end{proof}

\begin{lemma}[指数基的玻尔均值正交性]\label{lem:amplitude_exponential_mean}
对任意 $m,n \in \mathbb{Z}$，
\[
M\{e^{i(m-n)t}\}
= \lim_{T\to\infty}\frac{1}{2T}\int_{-T}^{T} e^{i(m-n)t}\,dt
= \delta_{mn}.
\]
\end{lemma}

\begin{proof}
当 $m=n$ 时，被积函数为 $1$，均值为 $1$。当 $m\neq n$ 时，
\[
\frac{1}{2T}\int_{-T}^{T} e^{i(m-n)t}\,dt
= \frac{\sin((m-n)T)}{(m-n)T} \to 0 \qquad (T\to\infty).
\]
故均值为 $\delta_{mn}$。
\end{proof}

\begin{proposition}[振幅级数的帕塞瓦尔恒等式]\label{prop:amplitude_parseval}
设 $f \in L^2([0,2\pi])$，其傅里叶系数为 $c_n=A_n e^{i\phi_n}$，振幅级数为 $\sum_n A_n e^{int}$。则
\[
\frac{1}{2\pi}\int_0^{2\pi} |f(t)|^2\,dt
= \sum_{n=-\infty}^{\infty} A_n^2.
\]
即帕塞瓦尔恒等式可以仅用振幅级数表述——能量守恒完全由振幅谱决定，与相位无关。
\end{proposition}

\begin{proof}
由帕塞瓦尔恒等式（定理 \ref{thm:parseval_fourier_classical}），
\[
\frac{1}{2\pi}\int_0^{2\pi} |f(t)|^2\,dt
= \sum_{n=-\infty}^{\infty} |c_n|^2
= \sum_{n=-\infty}^{\infty} A_n^2.
\]
其中 $A_n=|c_n|$。
\end{proof}

\begin{corollary}[能量仅由振幅决定]\label{cor:amplitude_energy_only}
设 $f$ 与 $g$ 有相同的振幅谱 $\{A_n\}$，但相位谱可不同。则
\[
\|f\|_{L^2([0,2\pi])} = \|g\|_{L^2([0,2\pi])}.
\]
因此 $L^2$ 能量是振幅级数的函数，与相位级数无关。
\end{corollary}

\begin{proof}
由命题 \ref{prop:amplitude_parseval}，两者能量均等于 $\sum_n A_n^2$。
\end{proof}

\begin{property}[振幅级数的基本性质]\label{prop:amplitude_series_properties}
\begin{enumerate}[label=(\roman*)]
    \item \textbf{非负性}：$A_n\ge 0$；
    \item \textbf{偶性}：若 $A_{-n}=A_n$，则实形式 $\sum_n A_n\cos(nt)$ 为偶函数；
    \item \textbf{平移不保持相位零}：时间平移 $t\mapsto t-\tau$ 一般会引入相位 $-n\tau$，故振幅级数本身不保留平移不变性；
    \item \textbf{能量可加性}：正交频率分量的能量可加，$\|\mathcal{A}_f\|^2=\sum_n A_n^2$；
    \item \textbf{衰减与光滑性}：$A_n$ 衰减越快，对应的周期函数越光滑。
\end{enumerate}
\end{property}

\begin{proof}
(i) 由定义。 (ii) 由 $\cos(-nt)=\cos(nt)$。 (iii) 平移后第 $n$ 项为 $A_n e^{in(t-\tau)}=A_n e^{-in\tau}e^{int}$，相位不再为零。 (iv) 由命题 \ref{prop:amplitude_parseval}。 (v) 逐项求导 $k$ 次产生因子 $(in)^k$，故 $A_n n^k$ 的衰减控制 $k$ 阶光滑性。
\end{proof}

\begin{example}[热核的振幅级数]\label{ex:amplitude_heat_kernel}
周期为 $2\pi$ 的热核
\[
H_t(x)=\sum_{n=-\infty}^{\infty} e^{-n^2 t} e^{inx}, \qquad t>0
\]
是相位为零、周期为 $2\pi$ 的振幅级数，其振幅为 $A_n=e^{-n^2 t}\ge 0$。由帕塞瓦尔恒等式，
\[
\frac{1}{2\pi}\int_0^{2\pi}|H_t(x)|^2\,dx
= \sum_{n=-\infty}^{\infty} e^{-2n^2 t}.
\]
当 $t\to 0^+$ 时，该级数在分布意义下趋于 $2\pi\sum_{k\in\mathbb{Z}}\delta(x-2\pi k)$。
\end{example}

\begin{proposition}[振幅级数的唯一性]\label{prop:amplitude_series_uniqueness}
设两个振幅级数 $\sum_n A_n e^{int}$ 与 $\sum_n B_n e^{int}$（$A_n, B_n \ge 0$）都一致绝对收敛，且作为函数相等。则 $A_n = B_n$ 对所有 $n$ 成立。
\end{proposition}

\begin{proof}
由玻尔均值公式（命题 \ref{prop:amplitude_bohr_mean}），$A_n = M\{\mathcal{A}_f e^{-int}\}$ 与 $B_n = M\{\mathcal{B}_f e^{-int}\}$。因 $\mathcal{A}_f = \mathcal{B}_f$（函数相等），故 $A_n = B_n$。
\end{proof}

\begin{example}[不同相位谱、同振幅谱的波形对比]\label{ex:same_amplitude_diff_phase}
取两个周期为 $2\pi$ 的信号：
\[
f_1(t) = \cos t + \frac{1}{2}\cos 2t, \qquad f_2(t) = \cos t + \frac{1}{2}\cos(2t + \pi/2) = \cos t - \frac{1}{2}\sin 2t.
\]
二者振幅谱相同（$A_1 = 1, A_2 = 1/2$），但相位谱不同（$f_2$ 的二次谐波相位为 $\pi/2$）。由推论 \ref{cor:amplitude_energy_only}，两者能量相同；但波形截然不同：$f_1$ 为偶函数且对称，$f_2$ 为非对称。
\end{example}

\begin{figure}[htbp]
\centering
\begin{tikzpicture}[scale=0.9,>=stealth]
  \draw[->] (-0.3,0) -- (7.0,0) node[right, font=\small] {$t$};
  \draw[->] (0,-1.7) -- (0,1.9) node[above, font=\small] {$f$};

  \draw[thick, blue!70!black, domain=0:6.6, samples=200]
    plot (\x, {cos(deg(\x)) + 0.5*cos(deg(2*\x))});
  \node[blue!70!black, font=\footnotesize] at (5.6,1.5) {$f_1(t)$};

  \draw[thick, red!70!black, dashed, domain=0:6.6, samples=200]
    plot (\x, {cos(deg(\x)) - 0.5*sin(deg(2*\x))});
  \node[red!70!black, font=\footnotesize] at (5.6,-1.4) {$f_2(t)$};

  \node[font=\small\bfseries] at (3.3,-2.5) {相同振幅谱，不同相位谱 $\Rightarrow$ 不同波形，相同能量};
\end{tikzpicture}
\caption{相同振幅谱、不同相位谱的两个信号对比。振幅谱相同保证能量相同（帕塞瓦尔，命题 \ref{prop:amplitude_parseval}），但波形因相位不同而显著不同。这说明了相位在波形重构中的不可替代作用。}
\label{fig:same_amplitude_diff_phase}
\end{figure}

\begin{remark}[振幅级数的适用范围与注意事项]\label{rem:amplitude_series_scope}
振幅级数只保留振幅 $A_n$，完全丢弃相位 $\phi_n$。因此它适用于能量估计、光滑性判断、衰减率分析等与相位无关的问题；但不能用于重构原函数、判断点态收敛位置、描述时间平移或波形对称性。若原级数不是绝对收敛，则“振幅级数”应理解为形式级数或分布级数；帕塞瓦尔恒等式需在 $L^2$ 或绝对收敛条件下使用。
\end{remark}

\section{\texorpdfstring{相位级数}{相位级数}}

本节讨论振幅为 $1$、周期为 $2\pi$ 的级数形式。它保留全部相位信息而丢弃振幅信息，是研究时间平移、频率调制与相位耦合的基本对象。

\begin{definition}[振幅为 $1$、周期为 $2\pi$ 的相位级数]\label{def:phase_series}
设 $\phi_n \in [0,2\pi)$（$n \in \mathbb{Z}$）为相位序列。形式级数
\[
\mathcal{P}_f(t) := \sum_{n=-\infty}^{\infty} e^{i(nt+\phi_n)}
= \sum_{n=-\infty}^{\infty} e^{i\phi_n} e^{int}
\]
称为\textbf{振幅为 $1$、周期为 $2\pi$ 的相位级数}。其生成函数形式为
\[
\mathfrak{P}_f(s) := \sum_{n=-\infty}^{\infty} e^{i\phi_n} s^n.
\]
相位级数的系数模恒为 $1$，其信息完全由辐角 $\phi_n$ 的分布决定。一般无限相位级数不逐点收敛，应按形式级数、有限截断或分布意义理解。
\end{definition}

\begin{proposition}[相位级数与时间平移]\label{prop:phase_translation}
设 $f(t)=\sum_n A_n e^{i(nt+\phi_n)}$ 的相位级数为 $\mathcal{P}_f(t)=\sum_n e^{i(nt+\phi_n)}$。则时间平移 $f_\tau(t):=f(t-\tau)$ 的相位级数为
\[
\mathcal{P}_{f_\tau}(t)
= \sum_n e^{i(nt+\phi_n-n\tau)}
= \sum_n e^{i(\phi_n-n\tau)}e^{int}.
\]
即时间平移等价于相位级数各项的旋转——第 $n$ 项旋转角 $-n\tau$ 与谐波序数 $n$ 成正比。
\end{proposition}

\begin{proof}
\[
f_\tau(t)=\sum_n A_n e^{i(n(t-\tau)+\phi_n)}
=\sum_n A_n e^{i(nt+\phi_n-n\tau)}.
\]
故相位分量为 $\phi_n-n\tau$。
\end{proof}

\begin{lemma}[相位级数的玻尔均值系数公式]\label{lem:phase_bohr_mean}
若相位级数 $\mathcal{P}_f(t)=\sum_n e^{i(nt+\phi_n)}$ 一致绝对收敛，则
\[
e^{i\phi_n}
= M\{\mathcal{P}_f(t)e^{-int}\}.
\]
\end{lemma}

\begin{proof}
由一致绝对收敛性，
\[
M\{\mathcal{P}_f(t)e^{-int}\}
= \sum_m e^{i\phi_m} M\{e^{i(m-n)t}\}
= e^{i\phi_n},
\]
其中最后一步用了引理 \ref{lem:amplitude_exponential_mean}。
\end{proof}

\begin{theorem}[有限相位级数的模长守恒]\label{thm:phase_modulus_conservation}
设 $S_N(t)=\sum_{n=-N}^{N} e^{i(nt+\phi_n)}$ 为有限相位级数。则
\[
\frac{1}{2\pi}\int_0^{2\pi}|S_N(t)|^2\,dt = 2N+1.
\]
即有限相位级数的总能量只由项数决定，与相位 $\phi_n$ 无关。
\end{theorem}

\begin{proof}
由正交性，
\[
|S_N(t)|^2
= \sum_{m,n=-N}^{N} e^{i((m-n)t+\phi_m-\phi_n)}.
\]
逐项积分：
\[
\frac{1}{2\pi}\int_0^{2\pi} e^{i(m-n)t}\,dt = \delta_{mn}.
\]
故仅对角项 $m=n$ 保留，每项贡献 $1$，总数为 $2N+1$。
\end{proof}

\begin{corollary}[时间平移的相位旋转规则]\label{cor:phase_time_shift_rotation}
时间平移 $t\mapsto t-\tau$ 将相位级数的第 $n$ 项乘以 $e^{-in\tau}$。因此：
\begin{enumerate}[label=(\roman*)]
    \item 若 $\tau=2\pi k$，则所有项不变；
    \item 若 $\tau=\pi$，则第 $n$ 项乘以 $(-1)^n$；
    \item 若 $\tau$ 为无理数倍 $2\pi$，则相位旋转不再具有有限阶周期性。
\end{enumerate}
\end{corollary}

\begin{proof}
由命题 \ref{prop:phase_translation}，第 $n$ 项相位变为 $\phi_n-n\tau$。逐条代入即得。
\end{proof}

\begin{property}[相位级数的基本性质]\label{prop:phase_series_properties}
\begin{enumerate}[label=(\roman*)]
    \item \textbf{单位模系数}：所有系数满足 $|e^{i\phi_n}|=1$；
    \item \textbf{共轭对称性}：若要求实信号，则需 $e^{i\phi_{-n}}=e^{-i\phi_n}$，即 $\phi_{-n}=-\phi_n \pmod{2\pi}$；
    \item \textbf{平移旋转性}：时间平移等价于逐项相位旋转；
    \item \textbf{能量项数守恒}：有限截断的总能量等于项数；
    \item \textbf{全无限一般不收敛}：$\sum_{n\in\mathbb{Z}}e^{int}$ 不逐点收敛，应在分布意义下理解。
\end{enumerate}
\end{property}

\begin{proof}
(i) 由定义。 (ii) 实信号要求 $c_{-n}=\overline{c_n}$，此处 $c_n=e^{i\phi_n}$，故 $e^{i\phi_{-n}}=e^{-i\phi_n}$。 (iii) 由命题 \ref{prop:phase_translation}。 (iv) 由定理 \ref{thm:phase_modulus_conservation}。 (v) 通项模为 $1$，不趋于零，故通常意义下发散；作为分布，$\sum_{n\in\mathbb{Z}}e^{int}=2\pi\sum_{k\in\mathbb{Z}}\delta(t-2\pi k)$。
\end{proof}

\begin{example}[线性相位与时间平移]\label{ex:linear_phase}
取 $\phi_n=-n\tau$，则
\[
\mathcal{P}(t)=\sum_{n=-N}^{N} e^{i(nt-n\tau)}
= \sum_{n=-N}^{N} e^{in(t-\tau)}
= D_N(t-\tau),
\]
其中 $D_N$ 为狄利克雷核。该例子说明线性相位对应时间平移。
\end{example}

\begin{proposition}[相位级数的旋转-平移对偶]\label{prop:phase_rotation_translation_duality}
设相位级数 $\mathcal{P}_f(t) = \sum_n e^{i(nt+\phi_n)}$。则：
\begin{enumerate}[label=(\roman*)]
  \item \textbf{时间平移} $t \mapsto t - \tau$ 等价于在频域中对第 $n$ 项乘以 $e^{-in\tau}$（命题 \ref{prop:phase_translation}）；
  \item \textbf{频率平移} $n \mapsto n + k$ 等价于在时域中对第 $n$ 项乘以 $e^{ikt}$。
\end{enumerate}
这是傅里叶变换的\textbf{平移-调制对偶}在相位级数上的体现。
\end{proposition}

\begin{proof}
(i) 由命题 \ref{prop:phase_translation} 直接给出。(ii) 频率平移 $n \mapsto n + k$ 将第 $n$ 项 $e^{i(nt+\phi_n)}$ 变为 $e^{i((n+k)t+\phi_{n+k})}$，即对原第 $n+k$ 项乘以 $e^{ikt}$。
\end{proof}

\begin{example}[线性相位与群延迟]\label{ex:linear_phase_group_delay}
设 $\phi_n = -n\tau$（线性相位）。则相位级数的群延迟为
\[
\tau_g(n) = -\frac{d\phi_n}{dn} = \tau,
\]
即所有频率分量的群延迟相同。在信号处理中，群延迟恒定的系统称为\textbf{线性相位系统}，对应时域中的纯平移 $t \mapsto t - \tau$（例 \ref{ex:linear_phase}）。
\end{example}

\begin{property}[相位级数的Parseval等式（有限项）]\label{prop:phase_parseval_finite}
设 $\mathcal{P}_N(t) = \sum_{n=-N}^{N} A_n e^{i(nt+\phi_n)}$ 为有限相位级数（$A_n > 0$ 未必为 $1$）。则
\[
\frac{1}{2\pi}\int_0^{2\pi} |\mathcal{P}_N(t)|^2\,dt = \sum_{n=-N}^{N} A_n^2.
\]
\end{property}

\begin{proof}
由傅里叶基的正交性（命题 \ref{prop:fourier_orthogonality}），交叉项积分为零，仅保留对角项。
\end{proof}

\begin{figure}[htbp]
\centering
\begin{tikzpicture}[scale=0.9,>=stealth]
  \begin{scope}
    \draw[->] (-0.3,0) -- (4.5,0) node[right, font=\small] {$n$};
    \draw[->] (0,-2.2) -- (0,2.2) node[above, font=\small] {$\phi_n$};

    \draw[thick, blue!70!black] (0.3,0) -- (4.2,-1.5);
    \node[blue!70!black, font=\footnotesize] at (2.5,1.6) {线性相位};
    \node[blue!70!black, font=\footnotesize] at (2.5,1.0) {$\phi_n = -n\tau$};
  \end{scope}

  \begin{scope}[xshift=6.0cm]
    \draw[->] (-0.3,0) -- (4.5,0) node[right, font=\small] {$n$};
    \draw[->] (0,-2.2) -- (0,2.2) node[above, font=\small] {$\phi_n$};

    \draw[thick, red!70!black, domain=0.3:4.2, samples=100]
      plot (\x, {1.5*sin(deg(2*\x))});
    \node[red!70!black, font=\footnotesize] at (2.5,1.9) {非线性相位};
    \node[red!70!black, font=\footnotesize] at (2.5,-1.8) {群延迟依赖于 $n$};
  \end{scope}
\end{tikzpicture}
\caption{相位谱的两种典型形态。左：线性相位（群延迟恒定），对应时域的纯平移。右：非线性相位（群延迟依赖于频率），对应时域的波形畸变。}
\label{fig:phase_spectrum_examples}
\end{figure}

\begin{remark}[相位级数的适用范围与注意事项]\label{rem:phase_series_scope}
相位级数只保留相位 $e^{i\phi_n}$，完全丢弃振幅 $A_n$。因此它适用于时间平移、相位耦合、频率调制、同步性分析等问题；但不能单独给出能量分布、收敛速率或函数值大小。全无限相位级数一般不收敛，必须配合有限截断、窗函数或分布理论使用。相位与频率的乘积项 $-n\tau$ 表明相位级数中的“旋转”与谐波序数 $n$ 绑定，这是频率—相位对偶的基本表现。
\end{remark}

\section{\texorpdfstring{频率级数与级数谱系}{频率级数与级数谱系}}

在振幅级数与相位级数之后，本节建立\textbf{频率级数}，并给出级数谱系的五层次。频率级数只保留频率格结构，完全丢弃振幅与相位信息。

\begin{definition}[频率级数]\label{def:frequency_series}
设 $f$ 以 $T$ 为周期，其傅里叶展开为
\[
f(t)=\sum_{n} A_n e^{i(\omega_n t+\phi_n)}, \qquad \omega_n=\frac{2\pi n}{T}.
\]
$f$ 的\textbf{频率级数}定义为仅保留频率分量的形式级数
\[
\mathcal{F}_f(s) := \sum_{n=-\infty}^{\infty} \omega_n s^n
= \frac{2\pi}{T}\sum_{n=-\infty}^{\infty} n s^n.
\]
频率级数丢弃全部振幅与相位信息，只保留各次谐波的频率格。频率级数完全由基本周期 $T$ 决定——所有频率都是基频 $\omega_0=2\pi/T$ 的整数倍。频率级数的缩放性质：周期变为 $\lambda T$（$\lambda>0$）使 $\omega_n\mapsto \omega_n/\lambda$，即频率级数整体缩放 $1/\lambda$。
\end{definition}

\begin{proposition}[频率级数与周期缩放]\label{prop:frequency_scaling}
设 $f$ 以 $T$ 为周期，频率级数为
\[
\mathcal{F}_f(s)=\frac{2\pi}{T}\sum_n n s^n.
\]
则时间缩放 $f_\lambda(t):=f(t/\lambda)$（$\lambda>0$）的频率级数为
\[
\mathcal{F}_{f_\lambda}(s)
= \frac{2\pi}{\lambda T}\sum_n n s^n
= \frac{1}{\lambda}\mathcal{F}_f(s).
\]
即时间缩放 $\lambda$ 倍等价于频率级数整体缩放 $1/\lambda$——这是频率-周期对偶在级数层面的表现。
\end{proposition}

\begin{proof}
$f_\lambda(t)=f(t/\lambda)$ 以 $\lambda T$ 为周期，故其傅里叶频率为
\[
\omega_n^{(\lambda)}=\frac{2\pi n}{\lambda T}=\frac{\omega_n}{\lambda}.
\]
故频率级数整体缩放 $1/\lambda$。
\end{proof}

\begin{definition}[级数谱系的五层次]\label{def:series_five_layers}
对以 $T$ 为周期的函数 $f$，其级数分解按信息量递增排列为：
\begin{enumerate}[label=(\roman*)]
    \item \textbf{频率级数} $\mathcal{F}_f$：仅保留频率格 $\omega_n=2\pi n/T$，完全由基本周期 $T$ 决定。
    \item \textbf{振幅级数} $\mathcal{A}_f$：保留振幅 $A_n$，丢弃相位 $\phi_n$ 与频率格结构。
    \item \textbf{相位级数} $\mathcal{P}_f$：保留相位 $e^{i\phi_n}$，丢弃振幅 $A_n$ 与频率格结构。
    \item \textbf{傅里叶级数}：保留振幅 $A_n$ 与相位 $\phi_n$，频率格由 $T$ 固定。
    \item \textbf{周期级数}：保留振幅 $A_{p,n}$、频率 $\omega_{p,n}$、相位 $\phi_{p,n}$ 与周期 $T_p$ 四元组，是多周期层的完整分解。
\end{enumerate}
五层次的关系为：傅里叶级数 $=$ 振幅级数 $\times$ 相位级数（逐项）；周期级数 $=$ 多周期层的傅里叶级数之和；振幅级数、相位级数、频率级数分别是傅里叶级数在振幅、相位、频率三个分量上的投影。
\end{definition}

\begin{proposition}[五层次的还原关系]\label{prop:series_five_layers_reconstruction}
在上述记号下：
\begin{enumerate}[label=(\roman*)]
    \item 由振幅级数与相位级数可还原傅里叶系数：$c_n=A_n e^{i\phi_n}$，从而还原傅里叶级数。
    \item 由振幅级数可还原能量：$\frac{1}{T}\int_0^T |f|^2=\sum_n A_n^2$（帕塞瓦尔）。
    \item 由相位级数可还原时间平移效应：$\phi_n-2\pi n\tau/T$。
    \item 由频率级数可还原周期缩放效应：$\omega_n\mapsto \omega_n/\lambda$。
    \item 五层次中任一单独层次的信息都不足以还原 $f$；四元组齐全时（周期级数）方可完整还原。
\end{enumerate}
\end{proposition}

\begin{proof}
(i) 由 $c_n=A_n e^{i\phi_n}$ 直接给出。
(ii) 由命题 \ref{prop:amplitude_parseval}。
(iii) 由命题 \ref{prop:phase_translation}。
(iv) 由命题 \ref{prop:frequency_scaling}。
(v) 由定义 \ref{def:series_five_layers} 的逐项信息对比即得：频率级数只含 $\omega_n$，振幅级数只含 $A_n$，相位级数只含 $e^{i\phi_n}$，任一都缺失其他分量的信息。
\end{proof}

\begin{example}[五种级数层次的显式对比]\label{ex:five_layers_example}
取周期 $T = 2\pi$ 的周期函数 $f(t) = 3 + 2 e^{i(t + \pi/3)} + e^{i(3t - \pi/4)}$。五种层次的级数分别为：
\begin{center}
\begin{tabular}{@{}p{2.8cm}p{9.0cm}@{}}
\toprule
层次 & 具体形式 \\
\midrule
频率级数 & $\mathcal{F}_f(s) = \sum_n n s^n$，仅保留频率格 $\{0, 1, 3\}$（模 $2\pi$） \\
振幅级数 & $\mathcal{A}_f(t) = 3 + 2 e^{it} + e^{i3t}$，振幅 $\{3, 2, 1\}$ \\
相位级数 & $\mathcal{P}_f(t) = e^{i\pi/3} e^{it} + e^{-i\pi/4} e^{i3t}$，相位 $\{\pi/3, -\pi/4\}$ \\
傅里叶级数 & $f(t) = 3 + 2 e^{i(t+\pi/3)} + e^{i(3t-\pi/4)}$ \\
周期级数 & 单周期层 $P = \{2\pi\}$ 上的傅里叶级数 \\
\bottomrule
\end{tabular}
\end{center}
该例子说明：五种层次中，只有傅里叶级数与周期级数（在单周期层情形下二者重合）可以完整还原 $f$；频率级数、振幅级数、相位级数分别只保留一个分量。
\end{example}

\begin{figure}[htbp]
\centering
\begin{tikzpicture}[scale=0.9,>=stealth,
  layer/.style={draw, rounded corners, minimum width=3.5cm, minimum height=1.0cm, align=center, font=\footnotesize},
  arrow/.style={->, thick, >=stealth}]

  \node[layer, fill=blue!15] (freq) at (0,0) {频率级数\\仅频率格 $\omega_n$};
  \node[layer, fill=green!15] (amp) at (0,-1.6) {振幅级数\\仅振幅 $A_n$};
  \node[layer, fill=orange!15] (phase) at (0,-3.2) {相位级数\\仅相位 $e^{i\phi_n}$};
  \node[layer, fill=red!15] (fourier) at (0,-4.8) {傅里叶级数\\振幅 $+$ 相位 $+$ 频率格};
  \node[layer, fill=purple!15] (periodic) at (0,-6.4) {周期级数\\振幅 $+$ 频率 $+$ 相位 $+$ 周期};

  \draw[arrow] (amp.east) to[bend left=20] node[right, font=\scriptsize, align=left] {投影} (freq.east);
  \draw[arrow] (phase.east) to[bend left=20] (amp.east);
  \draw[arrow] (fourier.east) to[bend left=20] (phase.east);
  \draw[arrow] (periodic.east) to[bend left=20] (fourier.east);

  \node[font=\small\bfseries, rotate=90] at (-3.5,-3.2) {信息量递增};

\end{tikzpicture}
\caption{级数谱系五层次。信息量从频率级数（仅频率格）到周期级数（四元组完整）逐层递增，每一层都是上一层的投影或综合。}
\label{fig:series_five_layers}
\end{figure}

\section{\texorpdfstring{周期级数：四元组的完整分解}{周期级数}}

\begin{definition}[周期级数（复形式）]\label{def:periodic_series}
设 $P$ 为有限或可数的\textbf{周期类指标集}，对每个 $p \in P$ 指定基本周期 $T_p > 0$，$n \in \mathbb{Z}$ 为该周期类内的谐波序数。形如
\[
f(t) = \sum_{p \in P} \sum_{n \in \mathbb{Z}} A_{p,n}\, e^{i(2\pi n t / T_p + \phi_{p,n})}
\]
的级数称为\textbf{周期级数}（复形式），其中第 $(p,n)$ 项由\textbf{四元组} $(A_{p,n}, \omega_{p,n}, \phi_{p,n}, T_p)$ 确定：
\begin{enumerate}[label=(\roman*)]
    \item $A_{p,n} \ge 0$ 称为\textbf{振幅分量}；
    \item $\omega_{p,n} := 2\pi n / T_p$ 称为\textbf{频率分量}；
    \item $\phi_{p,n} \in [0, 2\pi)$ 称为\textbf{相位分量}；
    \item $T_p$ 称为\textbf{周期分量}（一级指标）。
\end{enumerate}
频率与周期由\textbf{定长关系} $\omega_{p,n} = 2\pi n / T_p$ 绑定。记极式系数 $c_{p,n} = A_{p,n} e^{i\phi_{p,n}}$，则 $f(t) = \sum_{p,n} c_{p,n} e^{i\omega_{p,n} t}$。
\end{definition}

\begin{definition}[周期级数（实形式）]\label{def:periodic_series_real}
设 $f: \mathbb{R} \to \mathbb{R}$。形如
\[
f(t) = \sum_{p \in P} \sum_{n \in \mathbb{Z}} A_{p,n} \cos\!\big( \tfrac{2\pi n}{T_p} t + \phi_{p,n} \big), \qquad A_{p,n} \ge 0,\ \phi_{p,n} \in [0, 2\pi)
\]
的级数称为\textbf{周期级数}（实形式）。
\end{definition}

\begin{proposition}[实形式与复形式的等价]\label{prop:periodic_series_real_complex}
实形式周期级数是复形式周期级数的实部。
\end{proposition}

\begin{proof}
由欧拉公式 $e^{i\theta} = \cos\theta + i\sin\theta$ 直接得出。
\end{proof}

\begin{proposition}[傅里叶级数是周期级数的特殊情形]\label{prop:fourier_as_periodic_series}
傅里叶级数是周期级数在单周期层（$P = \{T\}$）上的特殊情形。
\end{proposition}

\begin{proof}
对 $T$-周期函数，玻尔均值化为单周期平均：
\[
M\{g\} = \frac{1}{T}\int_0^T g(t)\,dt.
\]
故周期级数退化为傅里叶级数。
\end{proof}

\begin{proposition}[周期级数对振幅级数、相位级数与频率级数的综合]\label{prop:periodic_series_synthesis}
设周期级数
\[
f(t)=\sum_{p\in P}\sum_{n\in\mathbb Z}A_{p,n}e^{i(\omega_{p,n}t+\phi_{p,n})},
\qquad \omega_{p,n}=\frac{2\pi n}{T_p}.
\]
对每个周期层 $p\in P$，记
\[
\mathcal A_p(s)=\sum_n A_{p,n}s^n,\qquad
\mathcal P_p(s)=\sum_n e^{i\phi_{p,n}}s^n,\qquad
\mathcal F_p(s)=\sum_n \omega_{p,n}s^n.
\]
则：
\begin{enumerate}[label=(\roman*)]
    \item 若固定单周期层 $P=\{T\}$、$T=2\pi$，且令所有 $\phi_{p,n}=0$，则 $f$ 退化为振幅级数
    \[
    \sum_n A_n e^{int};
    \]
    \item 若固定单周期层 $P=\{T\}$、$T=2\pi$，且令所有 $A_{p,n}=1$，则 $f$ 退化为相位级数
    \[
    \sum_n e^{i(nt+\phi_n)};
    \]
    \item 若固定单周期层并只保留频率参数 $\omega_n=2\pi n/T$，则得到频率级数；
    \item 一般周期级数由振幅 $A_{p,n}$、相位 $\phi_{p,n}$、频率 $\omega_{p,n}$ 与周期 $T_p$ 四元组按周期层 $p$ 综合而成。
\end{enumerate}
\end{proposition}

\begin{proof}
(i)(ii)(iii) 直接由定义 \ref{def:periodic_series}、\ref{def:amplitude_series}、\ref{def:phase_series}、\ref{def:frequency_series} 取相应参数即得。(iv) 由
\[
c_{p,n}=A_{p,n}e^{i\phi_{p,n}},\qquad \omega_{p,n}=\frac{2\pi n}{T_p}
\]
直接读出。
\end{proof}

\begin{corollary}[三种级数是周期级数的投影]\label{cor:periodic_series_projections}
振幅级数、相位级数与频率级数分别是周期级数在振幅、相位、频率三个分量上的投影。任意单一投影不足以重构 $f$；四元组齐全时可唯一重构。
\end{corollary}

\begin{proof}
由命题 \ref{prop:periodic_series_synthesis} 与命题 \ref{prop:series_five_layers_reconstruction}。
\end{proof}

\begin{remark}[“综合”的准确含义]\label{rem:periodic_series_synthesis_scope}
“周期级数是三种级数的综合”应理解为：以周期 $T_p$ 为一级指标，在每个周期层内由振幅级数提供 $A_{p,n}$、相位级数提供 $e^{i\phi_{p,n}}$、频率级数提供 $\omega_{p,n}=2\pi n/T_p$，再通过
\[
A_{p,n}e^{i(\omega_{p,n}t+\phi_{p,n})}
\]
组合成周期级数。它不是三个形式级数的逐项相加，而是四元组参数化的分层傅里叶综合。单周期层 $P=\{T\}$、$T=2\pi$ 时退化为傅里叶级数；进一步取相位为零或振幅为一分别得到振幅级数或相位级数。
\end{remark}

\begin{example}[准周期的周期级数]\label{ex:quasi_periodic_series}
\[
f(t) = 3 e^{i(\sqrt{2}\,t + \pi/6)} + 2 e^{i(t - \pi/3)}.
\]
频率 $\sqrt{2}$ 与 $1$ 不可公度，故 $f$ 是准周期函数。
\end{example}

\begin{property}[周期级数的傅里叶级数退化]\label{prop:periodic_series_degeneration}
当周期类指标集 $P = \{T\}$ 为单点集时，周期级数退化为以 $T$ 为周期的傅里叶级数：
\[
f(t) = \sum_{n \in \mathbb{Z}} A_n e^{i(2\pi n t / T + \phi_n)}.
\]
进一步，若 $T = 2\pi$，则退化为标准傅里叶级数 $f(t) = \sum_{n \in \mathbb{Z}} c_n e^{int}$。
\end{property}

\begin{proof}
单周期层 $P = \{T\}$ 时，外层求和消失，内层求和即傅里叶级数形式。取 $T = 2\pi$ 即得标准形式。
\end{proof}

\begin{example}[双周期层的周期级数]\label{ex:two_period_layers}
取 $P = \{2\pi, 2\pi\sqrt{2}\}$，构造
\[
f(t) = 3 e^{i(t + \pi/6)} + 2 e^{i(\sqrt{2}\,t - \pi/3)}.
\]
该函数有\textbf{两个独立的周期层}：外层为周期 $2\pi$ 的分量，内层为周期 $2\pi\sqrt{2}$ 的分量。由于 $2\pi/(2\pi\sqrt{2}) = 1/\sqrt{2} \notin \mathbb{Q}$，两层不可公度，$f$ 是准周期函数（例 \ref{ex:quasi_periodic_series}）。
\end{example}

\begin{theorem}[周期级数收敛的充分条件]\label{thm:periodic_series_convergence}
设周期级数 $f(t) = \sum_{p \in P} \sum_{n \in \mathbb{Z}} A_{p,n} e^{i(\omega_{p,n} t + \phi_{p,n})}$ 满足
\[
\sum_{p \in P} \sum_{n \in \mathbb{Z}} A_{p,n} < \infty.
\]
则 $f$ 在 $\mathbb{R}$ 上一致绝对收敛，且 $f$ 连续有界。
\end{theorem}

\begin{proof}
对任意 $t$，$|A_{p,n} e^{i(\omega_{p,n} t + \phi_{p,n})}| = A_{p,n}$。由魏尔斯特拉斯 M-判别法（定理 \ref{thm:weierstrass_mtest}），级数一致收敛。连续性与有界性由一致收敛保持。
\end{proof}

\subsection{\texorpdfstring{单位系数特例：相位为零、振幅为一}{单位系数特例：相位为零、振幅为一}}

本小节讨论周期级数在相位为零、振幅为一时的特殊形式。它是一般傅里叶级数的归一化基特例，也是理解正交展开与分布收敛的入口。

\begin{definition}[相位为零、振幅为一的周期级数]\label{def:unit_phase_zero_amplitude_series}
设 $S \subseteq \mathbb{Z}$ 为频率指标集，$T>0$。形式级数
\[
\Phi_{T,S}(t) := \sum_{n \in S} e^{i 2\pi n t/T}
\]
称为\textbf{相位为零、振幅为一的周期级数}。当 $T=2\pi$ 时，简记为
\[
\Phi_S(t) := \sum_{n \in S} e^{i n t}.
\]
其实形式（要求 $S=-S$）为
\[
\Phi_S^{\mathbb{R}}(t) := \sum_{n \in S} \cos(nt).
\]
若 $S$ 为有限集，则 $\Phi_{T,S}$ 为三角多项式，逐点绝对收敛；若 $S=\mathbb{Z}$，则应按分布意义理解。
\end{definition}

\begin{proposition}[系数、相位与振幅的归一化]\label{prop:unit_phase_zero_coefficients}
相位为零、振幅为一的周期级数 $\Phi_{T,S}(t)=\sum_{n\in S}e^{i2\pi n t/T}$ 的傅里叶系数为
\[
c_n =
\begin{cases}
1, & n\in S,\\
0, & n\notin S.
\end{cases}
\]
因此其振幅为 $A_n=1$（$n\in S$），相位为 $\phi_n=0$（$n\in S$）。
\end{proposition}

\begin{proof}
由傅里叶系数公式与正交性（命题 \ref{prop:fourier_orthogonality}），
\[
c_n=\frac{1}{T}\int_0^T \Phi_{T,S}(t)e^{-i2\pi n t/T}\,dt
=\sum_{m\in S}\frac{1}{T}\int_0^T e^{i2\pi (m-n)t/T}\,dt
=\sum_{m\in S}\delta_{mn}.
\]
故当 $n\in S$ 时 $c_n=1$，否则 $c_n=0$。
\end{proof}

\begin{lemma}[有限频率集的 $L^2$ 正交性]\label{lem:unit_phase_zero_orthogonality}
设 $S\subset \mathbb{Z}$ 有限。则
\[
\frac{1}{T}\int_0^T \Phi_{T,S}(t)\overline{\Phi_{T,S}(t)}\,dt = |S|,
\]
其中 $|S|$ 为 $S$ 中元素个数。更一般地，
\[
\frac{1}{T}\int_0^T e^{i2\pi m t/T}e^{-i2\pi n t/T}\,dt=\delta_{mn}.
\]
\end{lemma}

\begin{proof}
展开：
\[
\frac{1}{T}\int_0^T \Phi_{T,S}\overline{\Phi_{T,S}}\,dt
=\sum_{m,n\in S}\frac{1}{T}\int_0^T e^{i2\pi (m-n)t/T}\,dt
=\sum_{m,n\in S}\delta_{mn}=|S|.
\]
\end{proof}

\begin{theorem}[相位为零、振幅为一的周期级数与一般傅里叶级数的关系]\label{thm:unit_phase_zero_fourier_relation}
设 $f$ 以 $T$ 为周期，其傅里叶级数为
\[
f(t)\sim \sum_{n\in\mathbb{Z}} c_n e^{i2\pi n t/T},
\qquad c_n=A_n e^{i\phi_n}.
\]
则：
\begin{enumerate}[label=(\roman*)]
    \item 一般傅里叶级数是振幅级数与相位级数的逐项乘积：
    \[
    c_n e^{i2\pi n t/T}
    = A_n e^{i\phi_n}e^{i2\pi n t/T}.
    \]
    \item 相位为零、振幅为一的周期级数是傅里叶级数在 $c_n=1$（$n\in S$）时的归一化特例：
    \[
    \Phi_{T,S}(t)=\sum_{n\in S}1\cdot e^{i\phi_n}e^{i2\pi n t/T},
    \qquad \phi_n=0,\ A_n=1.
    \]
    \item 若 $S$ 有限，则 $\Phi_{T,S}$ 是 $L^2([0,T])$ 中由频率集 $S$ 张成的正交投影核；若 $S=\mathbb{Z}$，则 $\Phi_{T,\mathbb{Z}}$ 是周期狄利克雷核的分布极限：
    \[
    \sum_{n\in\mathbb{Z}}e^{i2\pi n t/T}
    = T\sum_{k\in\mathbb{Z}}\delta(t-kT).
    \]
\end{enumerate}
\end{theorem}

\begin{proof}
(i) 由 $c_n=A_n e^{i\phi_n}$ 直接写出。
(ii) 在 (i) 中取 $A_n=1$、$\phi_n=0$（$n\in S$），其余 $c_n=0$，即得。
(iii) 有限 $S$ 时，$\Phi_{T,S}$ 是有限个正交基函数的和，其积分核性质由引理 \ref{lem:unit_phase_zero_orthogonality} 给出。$S=\mathbb{Z}$ 时，利用泊松求和公式（或分布意义下的狄利克雷核极限）：
\[
\sum_{n\in\mathbb{Z}}e^{i2\pi n t/T}
= T\sum_{k\in\mathbb{Z}}\delta(t-kT).
\]
故应按分布理解。
\end{proof}

\begin{corollary}[有限归一化基的帕塞瓦尔恒等式]\label{cor:unit_phase_zero_parseval}
设 $S\subset\mathbb{Z}$ 有限。则
\[
\frac{1}{T}\int_0^T |\Phi_{T,S}(t)|^2\,dt = |S|.
\]
即相位为零、振幅为一的有限周期级数的能量等于频率项数。
\end{corollary}

\begin{proof}
由引理 \ref{lem:unit_phase_zero_orthogonality} 直接得到。
\end{proof}

\begin{property}[相位为零、振幅为一的周期级数的基本性质]\label{prop:unit_phase_zero_properties}
\begin{enumerate}[label=(\roman*)]
    \item \textbf{系数单位性}：在 $S$ 上 $c_n=1$，振幅为 $1$，相位为 $0$；
    \item \textbf{正交性}：不同频率项在 $L^2$ 中正交；
    \item \textbf{有限能量可加}：有限截断能量等于项数；
    \item \textbf{实形式为余弦和}：若 $S=-S$，则实形式为 $\sum_{n\in S}\cos(2\pi n t/T)$；
    \item \textbf{全频率分布性}：$S=\mathbb{Z}$ 时不逐点收敛，而是周期 $\delta$ 梳。
\end{enumerate}
\end{property}

\begin{proof}
(i) 由命题 \ref{prop:unit_phase_zero_coefficients}。 (ii) 由引理 \ref{lem:unit_phase_zero_orthogonality}。 (iii) 由推论 \ref{cor:unit_phase_zero_parseval}。 (iv) 由欧拉公式。 (v) 由定理 \ref{thm:unit_phase_zero_fourier_relation}(iii)。
\end{proof}

\begin{example}[狄利克雷核作为有限归一化周期级数]\label{ex:dirichlet_kernel_truncation}
取 $T=2\pi$，$S=\{-N,\dots,N\}$，则
\[
\Phi_{\{-N,\dots,N\}}(t)
=\sum_{n=-N}^{N}e^{int}
=D_N(t)
=\frac{\sin((N+\tfrac12)t)}{\sin(t/2)}.
\]
其能量为
\[
\frac{1}{2\pi}\int_0^{2\pi}|D_N(t)|^2\,dt=2N+1.
\]
当 $N\to\infty$ 时，$D_N$ 在分布意义下趋于 $2\pi\sum_{k\in\mathbb{Z}}\delta(t-2\pi k)$。
\end{example}

\begin{remark}[相位为零、振幅为一的周期级数的适用范围]\label{rem:unit_phase_zero_scope}
相位为零、振幅为一的周期级数是傅里叶级数的“单位系数”特例，适用于说明正交基、狄利克雷核、泊松求和与分布收敛。它不表示一般 $L^2$ 或连续周期函数，因为全频率单位系数级数不收敛；只有有限频率集或分布意义下才有严格意义。对实信号，还需频率集关于原点对称并取实部。与一般傅里叶级数的关系是：一般级数 $=$ 振幅级数 $\times$ 相位级数，而本小节情形是振幅全为 $1$、相位全为 $0$ 的归一化交集。
\end{remark}

\subsection{\texorpdfstring{玻尔均值与系数的确定}{玻尔均值与系数的确定}}

\begin{definition}[玻尔均值]\label{def:bohr_mean}
设 $f: \mathbb{R} \to \mathbb{C}$。若极限
\[
M\{f\} := \lim_{T \to \infty} \frac{1}{2T} \int_{-T}^{T} f(t)\,dt
\]
存在，则称其为 $f$ 的\textbf{玻尔均值}。
\end{definition}

\begin{lemma}[指数函数的玻尔均值]\label{lem:bohr_mean_exponential}
对任意 $\omega \in \mathbb{R}$，
\[
M\{e^{i\omega t}\} =
\begin{cases}
1, & \omega = 0, \\
0, & \omega \neq 0.
\end{cases}
\]
\end{lemma}

\begin{proof}
$\omega = 0$ 时被积函数恒为 $1$。$\omega \neq 0$ 时，
\[
\frac{1}{2T}\int_{-T}^{T} e^{i\omega t}\,dt = \frac{\sin(\omega T)}{\omega T} \to 0 \qquad (T \to \infty).
\]
\end{proof}

\begin{proposition}[有限情形的初等证明：范德蒙行列式]\label{prop:vandermonde_uniqueness}
设频率 $\omega_1 < \cdots < \omega_N$ 两两互异。若 $\sum_{k=1}^{N} c_k e^{i\omega_k t} \equiv 0$，则 $c_1 = \cdots = c_N = 0$。
\end{proposition}

\begin{proof}
对恒等式逐次求导 $j = 0, \dots, N-1$ 次并在 $t = 0$ 取值，得齐次线性方程组，系数矩阵为范德蒙矩阵，行列式非零。
\end{proof}

\begin{theorem}[周期级数系数的确定（玻尔）]\label{thm:periodic_series_bohr}
设周期级数 $f(t) = \sum_{p,n} A_{p,n} e^{i(\omega_{p,n} t + \phi_{p,n})}$ 在 $\mathbb{R}$ 上一致收敛，且全部频率两两互异。则
\[
A_{p,n} e^{i\phi_{p,n}} = M\big\{ f(t)\, e^{-i\omega_{p,n} t} \big\}.
\]
\end{theorem}

\begin{proof}
由一致收敛性，积分与求和可交换。由引理 \ref{lem:bohr_mean_exponential}，交叉项均值为零，仅 $(p,n)$ 项保留。
\end{proof}

\begin{corollary}[表示的唯一性]\label{cor:periodic_series_unique}
若两个一致收敛的周期级数作为函数相等，且周期划分 $P$ 与频率集合相同，则两者的振幅与相位逐项相同。
\end{corollary}

\begin{proof}
由定理 \ref{thm:periodic_series_bohr} 直接得出。
\end{proof}

\begin{theorem}[广义帕塞瓦尔恒等式（玻尔）]\label{thm:bohr_parseval}
设周期级数 $f(t) = \sum_{p,n} c_{p,n} e^{i\omega_{p,n} t}$ 满足 $\sum_{p,n} |c_{p,n}| < \infty$。则
\[
M\{|f|^2\} = \sum_{p,n} |c_{p,n}|^2 = \sum_{p,n} A_{p,n}^2.
\]
\end{theorem}

\begin{proof}
展开 $|f(t)|^2$，由绝对收敛性玻尔均值可与求和交换，交叉项均值为零，对角项给出 $|c_{p,n}|^2$。
\end{proof}

\subsection{\texorpdfstring{周期变换：经典积分变换的统一}{周期变换}}

\begin{definition}[周期变换（频率字典形式）]\label{def:periodic_transform}
设 $\Lambda = \{\omega_k\}_{k \in K}$ 为\textbf{频率字典}。对函数 $f$，其关于字典 $\Lambda$ 的\textbf{周期变换}是系数族
\[
\mathcal{P}_\Lambda f(k) = M\big\{ f(t)\, e^{-i\omega_k t} \big\}, \qquad k \in K,
\]
其合成（反变换）为周期级数
\[
f(t) \sim \sum_{k \in K} \mathcal{P}_\Lambda f(k)\, e^{i\omega_k t}.
\]
\end{definition}

\begin{proposition}[周期变换的线性性]\label{prop:periodic_transform_linear}
设 $f, g$ 的周期变换关于字典 $\Lambda = \{\omega_k\}_{k \in K}$ 均存在，$\alpha, \beta \in \mathbb{C}$。则对每个 $k \in K$，
\[
\mathcal{P}_\Lambda(\alpha f + \beta g)(k) = \alpha\, \mathcal{P}_\Lambda f(k) + \beta\, \mathcal{P}_\Lambda g(k),
\]
即 $\mathcal{P}_\Lambda$ 是从函数空间到系数族空间 $\mathbb{C}^K$ 的\textbf{线性映射}。
\end{proposition}

\begin{proof}
分三步。

\textbf{第一步：$M\{\cdot\}$ 的线性}。设 $M\{f\}$ 与 $M\{g\}$ 均存在。由定义 \ref{def:bohr_mean}，$M\{\cdot\}$ 是部分积分平均的极限：
\[
M\{\alpha f + \beta g\}
= \lim_{T \to \infty} \frac{1}{2T}\int_{-T}^{T} \bigl( \alpha f(t) + \beta g(t) \bigr)\,dt .
\]
由于定积分关于被积函数线性，
\[
\frac{1}{2T}\int_{-T}^{T} \bigl( \alpha f(t) + \beta g(t) \bigr)\,dt
= \alpha \cdot \frac{1}{2T}\int_{-T}^{T} f(t)\,dt + \beta \cdot \frac{1}{2T}\int_{-T}^{T} g(t)\,dt .
\]
右端两个极限均存在（分别为 $M\{f\}$ 与 $M\{g\}$），故由极限的加法与数乘法则，
\[
M\{\alpha f + \beta g\} = \alpha M\{f\} + \beta M\{g\} .
\]
\textbf{这要求两个均值分别存在}：若其中之一不存在，右端无意义，线性性亦无从谈起。

\textbf{第二步：逐频率代入}。固定 $k \in K$，记
\[
f_k(t) := f(t)\,e^{-i\omega_k t},\qquad g_k(t) := g(t)\,e^{-i\omega_k t}.
\]
由 $M\{f_k\} = \mathcal{P}_\Lambda f(k)$ 与 $M\{g_k\} = \mathcal{P}_\Lambda g(k)$（定义 \ref{def:periodic_transform}），且
\[
(\alpha f + \beta g)(t)\,e^{-i\omega_k t} = \alpha f_k(t) + \beta g_k(t),
\]
第一步给出
\[
\mathcal{P}_\Lambda(\alpha f + \beta g)(k)
= M\{\alpha f_k + \beta g_k\}
= \alpha M\{f_k\} + \beta M\{g_k\}
= \alpha\, \mathcal{P}_\Lambda f(k) + \beta\, \mathcal{P}_\Lambda g(k).
\]

\textbf{第三步：逐点线性即映射线性}。上式对每个 $k \in K$ 成立，故两个系数族作为 $\mathbb{C}^K$ 中的元素满足线性关系；按线性映射的定义，$\mathcal{P}_\Lambda$ 线性。证毕。
\end{proof}

\begin{remark}[线性性的边界]
命题 \ref{prop:periodic_transform_linear} 只用到玻尔均值的线性，与频率字典 $\Lambda$ 的具体取法无关。但须注意两点：其一，线性性以各均值\textbf{分别存在}为前提，$M\{\cdot\}$ 本身并非在整个函数空间上有定义；其二，合成（反变换）$f \sim \sum_k \mathcal{P}_\Lambda f(k) e^{i\omega_k t}$ 的收敛性是另一个问题，线性性并不蕴含合成收敛。当 $\Lambda$ 取遍全部频率时，$\mathcal{P}_\Lambda$ 即经典傅里叶系数的求取算子（定理 \ref{thm:classical_transforms_periodic}）。
\end{remark}

\begin{theorem}[经典积分变换作为周期变换的特殊情形]\label{thm:classical_transforms_periodic}
下列经典变换都是周期变换在特定频率字典（或特定核）下的特殊情形：
\begin{enumerate}[label=(\roman*)]
    \item \textbf{傅里叶级数}：字典 $\Lambda = \omega_0\mathbb{Z}$。
    \item \textbf{离散傅里叶变换（DFT）}：有限循环群 $C_N$ 上，字典 $\Lambda = \{\frac{2\pi k}{N}\}_{k=0}^{N-1}$。
    \item \textbf{梅林变换}：经对数换元 $t = \ln x$：
    \[
    \mathcal{M}_f(s) = \int_0^\infty f(x)\, x^{s-1}\,dx = \int_{-\infty}^{\infty} f(e^t)\, e^{st}\,dt.
    \]
    \item \textbf{拉普拉斯变换}：复频率字典 $s = \sigma + i\omega$。
    \item \textbf{小波变换}：核化周期变换。
    \item \textbf{Chirp-Z 变换}：螺旋采样点列。
\end{enumerate}
\end{theorem}

\begin{proof}
(i)(ii) 代入周期变换定义即得。(iii) 换元 $x = e^t$ 直接给出。(iv) 复频率对应衰减因子。(v)(vi) 由各自定义直接验证。
\end{proof}

\begin{remark}[离散谱与连续谱]\label{rem:spectrum_decomposition}
玻尔均值只对具有离散谱的信号非零。一般信号的谱按勒贝格分解为原子部分（周期级数）、绝对连续部分（傅里叶变换密度）、奇异连续部分（分形/标度周期信号）。
\end{remark}

\section{\texorpdfstring{卷积：时延乘积与相位交互的统一算子}{卷积}}

\begin{definition}[卷积（时延乘积）]\label{def:convolution}
设 $f, g: \mathbb{R} \to \mathbb{R}$（或 $\mathbb{C}$）。$f$ 与 $g$ 的\textbf{卷积}是时延 $\tau$ 的函数：
\[
(f * g)(\tau) := \int_{-\infty}^{\infty} f(t) g(\tau - t)\,dt.
\]
当 $f, g$ 以 $T$ 为周期时，采用周期形式：
\[
(f * g)(\tau) := \frac{1}{T}\int_0^T f(t) g(\tau - t)\,dt.
\]
\end{definition}

\begin{definition}[互相关（共轭时延内积）]\label{def:cross_correlation}
设 $f, g$ 以 $T$ 为周期。$f$ 与 $g$ 的\textbf{互相关}定义为
\[
(f \star g)(\tau) := \frac{1}{T}\int_0^T f(t)\, \overline{g(t - \tau)}\,dt.
\]
\end{definition}

\begin{proposition}[卷积的基本代数性质]\label{prop:convolution_properties}
\begin{enumerate}[label=(\roman*)]
    \item \textbf{交换律}：$f * g = g * f$。
    \item \textbf{结合律}：$(f * g) * h = f * (g * h)$。
    \item \textbf{分配律}：$f * (g + h) = f * g + f * h$。
    \item \textbf{平移性质}：$\tau_a(f * g) = (\tau_a f) * g = f * (\tau_a g)$。
    \item \textbf{微分性质}：$(f * g)' = f' * g = f * g'$。
\end{enumerate}
\end{proposition}

\begin{proof}
交换律由变量代换 $u = \tau - t$ 得出；结合律与分配律由富比尼定理交换积分次序得出；平移性质由定义直接得出；微分性质由积分号下求导得出。
\end{proof}

\begin{proposition}[互相关与卷积的关系及频域形式]\label{prop:cross_corr_properties}
\begin{enumerate}[label=(\roman*)]
    \item 定义 $\widetilde{g}(t) = \overline{g(-t)}$，则 $(f \star g)(\tau) = (f * \widetilde{g})(\tau)$。
    \item $(f \star g)(\tau) = \overline{(g \star f)(-\tau)}$。
    \item 频域形式：$(f * g)(\tau) = \sum_{n} c_n d_n e^{i 2\pi n \tau / T}$，$(f \star g)(\tau) = \sum_{n} c_n \overline{d_n} e^{i 2\pi n \tau / T}$。
\end{enumerate}
\end{proposition}

\begin{proof}
(i) 直接代入定义。(ii) 取共轭并换元。(iii) 代入周期级数并利用正交性（命题 \ref{prop:fourier_orthogonality}）。
\end{proof}

\begin{example}[时延的物理意义：匹配滤波]\label{ex:echo_delay}
雷达发射信号 $s(t)$，接收回波 $r(t) = \alpha s(t - \tau_0) + \text{噪声}$。互相关 $(r \star s)(\tau)$ 在 $\tau = \tau_0$ 处取得峰值。
\end{example}

\begin{proposition}[卷积与傅里叶变换的乘积定理]\label{prop:convolution_fourier_product}
设 $f, g$ 以 $T$ 为周期，傅里叶系数分别为 $c_n, d_n$。则卷积 $f * g$ 的傅里叶系数为 $T c_n d_n$。
\end{proposition}

\begin{proof}
\[
\widehat{f * g}(n) = \frac{1}{T}\int_0^T (f * g)(t) e^{-i\omega_n t}\,dt = \frac{1}{T^2}\int_0^T \int_0^T f(\tau) g(t-\tau) e^{-i\omega_n t}\,d\tau\,dt.
\]
换元 $u = t - \tau$，利用富比尼定理：
\[
= \frac{1}{T} \int_0^T f(\tau) e^{-i\omega_n \tau}\,d\tau \cdot \int_0^T g(u) e^{-i\omega_n u}\,du = T c_n d_n.
\]
\end{proof}

\begin{corollary}[卷积算子的对角化]\label{cor:convolution_diagonalisation}
卷积算子 $C_f: g \mapsto f * g$ 在傅里叶基 $\{e^{i\omega_n t}\}_{n \in \mathbb{Z}}$ 下对角化：$C_f(e^{i\omega_n t}) = T c_n e^{i\omega_n t}$。即傅里叶基是卷积算子的特征函数，特征值为 $T c_n$。
\end{corollary}

\begin{proof}
由命题 \ref{prop:convolution_fourier_product}，$C_f(e^{i\omega_n t})$ 的傅里叶系数仅在 $n$ 处为 $T c_n$，其余为零，故 $C_f(e^{i\omega_n t}) = T c_n e^{i\omega_n t}$。
\end{proof}

\begin{example}[卷积的三个代数性质]\label{ex:convolution_algebra_example}
\begin{enumerate}[label=(\roman*)]
  \item \textbf{交换律}：$f * g = g * f$。例如 $e^{it} * e^{i2t} = e^{i2t} * e^{it}$（均为零，因为非同频率）；
  \item \textbf{结合律}：$(f * g) * h = f * (g * h)$，由卷积定理乘积结合律保证；
  \item \textbf{分配律}：$f * (g + h) = f * g + f * h$，由积分线性性保证。
\end{enumerate}
\end{example}

\begin{remark}[卷积与乘积的傅里叶对偶]\label{rem:convolution_product_duality}
卷积与乘积通过傅里叶变换构成一对\textbf{对偶}：
\[
\widehat{f * g} = \widehat{f} \cdot \widehat{g}, \qquad \widehat{f \cdot g} = \widehat{f} * \widehat{g}.
\]
前者是命题 \ref{prop:convolution_fourier_product}，后者由对偶性直接给出。这一对偶在信号处理中表现为：时域卷积等于频域乘积，时域乘积等于频域卷积（在适当的归一化下）。周期卷积（定义 \ref{def:periodic_convolution}）正是这一对偶在“保留全部频率对交互”这一目标下的推广。
\end{remark}

\begin{figure}[htbp]
\centering
\begin{tikzpicture}[>=stealth,
  box/.style={draw, rounded corners, align=center, font=\small, inner sep=6pt}]
  \node[box, fill=blue!8, minimum width=2.9cm] (f) at (-5,1.5) {$f$：$\{A_n,\omega_n,\phi_n\}$};
  \node[box, fill=orange!12, minimum width=3.6cm] (C) at (0,1.5) {交互张量\\[2pt]$C_{n,m}=A_nB_m\operatorname{Similarity}^{(d)}(\theta_n,\theta_m)$};
  \node[box, fill=blue!8, minimum width=2.9cm] (g) at (5,1.5) {$g$：$\{B_m,\nu_m,\psi_m\}$};
  
  \node[box, fill=green!12, minimum width=3.9cm] (S) at (0,-1.5) {合成映射\\[2pt]$S(C)(x)=\sum_{n,m}C_{n,m}e^{i(\Omega_{n,m}x+\Phi_{n,m})}$};
  
  \node[box, fill=red!8, minimum width=2.6cm] (out) at (0,-4.5) {$f\circledast g$};
  
  \draw[->,thick] (f) -- (C);
  \draw[->,thick] (g) -- (C);
  \draw[->,thick] (C) -- (S);
  \draw[->,thick] (S) -- (out);
  
  \node[font=\footnotesize] at (0,2.7) {第一层：分量两两交互};
  \node[font=\footnotesize] at (0,-0.2) {第二层：回到函数空间};
\end{tikzpicture}
\caption{卷积算子的两层构造。第一层把两个周转圆族的分量两两配对，得到双指标对象（$n\neq m$ 处即跨频率交互）；第二层由频率合成算子 $\Omega_{n,m}$ 与相位合成算子 $\Phi_{n,m}$ 把该对象送回函数空间。两层复合即卷积算子；当 $\operatorname{Similarity}^{(d)}$ 取相位相似核、$\Omega,\Phi$ 取和角形式时，即得到周期卷积 $f\circledast g$。}
\label{fig:convolution_two_layers}
\end{figure}

\subsection{\texorpdfstring{振幅卷积、相位卷积与频率卷积}{振幅卷积、相位卷积与频率卷积}}

设 $f, g$ 以 $T$ 为周期，其频率级数展开为
\[
f(t) = \sum_{m} c_m e^{i 2\pi m t/T}, \quad c_m = A_m e^{i\theta_m}; \qquad
g(t) = \sum_{n} d_n e^{i 2\pi n t/T}, \quad d_n = B_n e^{i\varphi_n}.
\]
其中 $A_m,B_n\ge 0$ 为振幅，$\theta_m,\varphi_n\in[0,2\pi)$ 为相位，$\omega_m=2\pi m/T$、$\nu_n=2\pi n/T$ 为频率。

\begin{definition}[振幅卷积]\label{def:amplitude_convolution}
$f$ 与 $g$ 的\textbf{振幅卷积}是实数矩阵 $\mathbb{A} = (\mathbb{A}_{mn})$：
\[
\mathbb{A}_{mn} := A_m B_n.
\]
振幅卷积丢弃全部相位信息，只保留尺度关系。$\mathbb{A}_{mn} \ge 0$ 对所有 $m, n$ 成立。
\end{definition}

\begin{definition}[相位卷积]\label{def:phase_convolution}
$f$ 与 $g$ 的\textbf{相位卷积}是复数矩阵 $\mathbb{P} = (\mathbb{P}_{mn})$：
\[
\mathbb{P}_{mn} := e^{i(\theta_m - \varphi_n)} = \cos(\theta_m - \varphi_n) + i\sin(\theta_m - \varphi_n).
\]
相位卷积丢弃全部振幅信息，只保留方向关系。$|\mathbb{P}_{mn}| = 1$ 对所有 $m, n$ 成立。
\end{definition}

\begin{definition}[频率卷积]\label{def:frequency_convolution}
$f$ 与 $g$ 的\textbf{频率卷积}是实数矩阵 $\mathbb{F} = (\mathbb{F}_{mn})$：
\[
\mathbb{F}_{mn} := \omega_m \nu_n = \frac{2\pi m}{T} \cdot \frac{2\pi n}{T} = \frac{4\pi^2 mn}{T^2}.
\]
频率卷积丢弃全部振幅与相位信息，只保留频率格结构。频率卷积完全由基本周期 $T$ 决定。
\end{definition}

\begin{definition}[趋势卷积]\label{def:tendency_convolution}
$f$ 与 $g$ 的\textbf{趋势卷积}是复数矩阵 $\mathbb{T} = (\mathbb{T}_{mn})$：
\[
\mathbb{T}_{mn} := A_m B_n \, \mathbb{P}_{mn} = A_m B_n e^{i(\theta_m - \varphi_n)} = c_m \overline{d_n}.
\]
趋势卷积 $=$ 振幅卷积 $\times$ 相位卷积（逐项）。
\end{definition}

\begin{proposition}[三类卷积与趋势卷积的关系]\label{prop:three_convolutions_relation}
在上述记号下：
\begin{enumerate}[label=(\roman*)]
    \item 趋势卷积 $=$ 振幅卷积 $\times$ 相位卷积：$\mathbb{T}_{mn} = \mathbb{A}_{mn} \mathbb{P}_{mn}$。
    \item 频率卷积 $=$ 频率级数外积：$\mathbb{F}_{mn} = \omega_m \nu_n$。
    \item 三者在频域中分别捕获卷积的三个独立分量：振幅（缩放）、相位（旋转）、频率（平移）。
\end{enumerate}
\end{proposition}

\begin{proof}
(i) 由定义 \ref{def:amplitude_convolution} 与 \ref{def:phase_convolution}：$\mathbb{A}_{mn}\mathbb{P}_{mn} = A_m B_n e^{i(\theta_m - \varphi_n)} = \mathbb{T}_{mn}$。
(ii) 由定义 \ref{def:frequency_convolution}。
(iii) 由 (i)(ii) 与三者的定义直接得出。
\end{proof}

\begin{proposition}[相位卷积的共轭对称性]\label{prop:phase_convolution_symmetry}
相位卷积矩阵满足 $\mathbb{P}_{mn} = \overline{\mathbb{P}'_{nm}}$，其中 $\mathbb{P}'_{mn} := e^{i(\varphi_m - \theta_n)}$ 为交换 $f$ 与 $g$ 的角色后所得的相位卷积。振幅卷积与频率卷积是对称实矩阵：$\mathbb{A}_{mn} = \mathbb{A}_{nm}$，$\mathbb{F}_{mn} = \mathbb{F}_{nm}$。
\end{proposition}

\begin{proof}
$\overline{\mathbb{P}_{nm}} = \overline{e^{i(\theta_n - \varphi_m)}} = e^{i(\varphi_m - \theta_n)}$，而 $\mathbb{P}_{mn} = e^{i(\theta_m - \varphi_n)}$；一般情形下二者不相等。正确的对称性为：定义 $\mathbb{P}'_{mn} := e^{i(\varphi_m - \theta_n)}$，则 $\mathbb{P}'_{mn} = \overline{\mathbb{P}_{nm}}$。振幅卷积与频率卷积的对称性由 $A_m B_n = B_n A_m$ 与 $\omega_m \nu_n = \nu_n \omega_m$ 直接得出。
\end{proof}

\begin{remark}[相位卷积的非对称性]\label{rem:phase_convolution_asymmetry}
相位卷积矩阵一般不是对称的：$\mathbb{P}_{mn} = e^{i(\theta_m - \varphi_n)}$ 与 $\mathbb{P}_{nm} = e^{i(\theta_n - \varphi_m)}$ 一般不相等，因为 $f$ 与 $g$ 的角色不可互换——$\theta_m$ 来自 $f$，$\varphi_n$ 来自 $g$。只有取复共轭后才有 $\mathbb{P}_{mn} = \overline{\mathbb{P}'_{nm}}$ 的联系，其中 $\mathbb{P}'_{mn} := e^{i(\varphi_m - \theta_n)}$ 为交换 $f$ 与 $g$ 的角色后所得的相位卷积。这与振幅卷积、频率卷积的对称性形成对比。
\end{remark}

\begin{definition}[核化卷积（核分解主公式）]\label{def:kernelized_convolution}
设权重核 $W: \mathbb{R} \to \mathbb{C}$、相似度核 $S^{(d)}$（$d = 1,\dots,6$）与振幅开关 $\alpha \in \{0, 1\}$。定义\textbf{核化卷积}
\[
\mathcal{K}^{(d)}_{W,\alpha}(\tau) := \sum_{m,n} W(\omega_m - \omega_n) \cdot (A_m B_n)^{\alpha} \cdot S^{(d)}(\theta_m, \varphi_n) \cdot e^{i\omega_n \tau}.
\]
\begin{enumerate}[label=(\roman*)]
    \item $\alpha = 0$：相位卷积（取 $S^{(d)}$ 为 $d$ 维相位核）；
    \item $\alpha = 1$：趋势卷积。
\end{enumerate}
\end{definition}

\begin{proposition}[6 个复相位核]\label{prop:six_complex_kernels}
6 个复相位核由差角/和角、基础/倍角/半角生成：
\[
S^{(1)} = e^{i(\theta_m - \varphi_n)}, \quad S^{(2)} = e^{i(\theta_m + \varphi_n)}, \quad S^{(3)} = e^{i \cdot 2(\theta_m - \varphi_n)},
\]
\[
S^{(4)} = e^{i(\theta_m - \varphi_n)/2}, \quad S^{(5)} = e^{i \cdot 2(\theta_m + \varphi_n)}, \quad S^{(6)} = e^{i(\theta_m + \varphi_n)/2}.
\]
\end{proposition}

\begin{proof}
“标量与矢量”章的 12 种实相似度恰为这 6 个复核的实部与虚部（余弦族与正弦族）。
\end{proof}

\begin{lemma}[带窗展开]\label{lem:windowed_expansion}
设 $w$ 为 $[-T, T]$ 上的可积窗，$\widehat{w}(\Omega) := \frac{1}{2T} \int_{-T}^{T} w(t) e^{i\Omega t}\,dt$。则
\[
\frac{1}{2T}\int_{-T}^{T} f(t) e^{-i\Omega t} w(t)\,dt = \sum_{m} c_m\, \widehat{w}(\omega_m - \Omega).
\]
\end{lemma}

\begin{proof}
代入 $f$ 的频率级数，由一致收敛性逐项积分。
\end{proof}

\begin{theorem}[窗-核对偶定理]\label{thm:window_kernel_duality}
设 $f, g$ 以 $T$ 为周期，$w$ 为偶窗。带窗时延内积
\[
(f \star_w g)(\tau) := \frac{1}{2T}\int_{-T}^{T} f(t)\, \overline{g(t - \tau)}\, w(t)\,dt
\]
具有展开
\[
(f \star_w g)(\tau) = \sum_{m,n} \widehat{w}(\omega_m - \omega_n) \cdot c_m \overline{d_n} \cdot e^{i\omega_n \tau}.
\]
\end{theorem}

\begin{proof}
代入频率级数，由引理 \ref{lem:windowed_expansion} 即得。
\end{proof}

\begin{corollary}[高斯窗与高斯核]\label{cor:gaussian_window_kernel}
取高斯时间窗 $w(t) = e^{-t^2/(2\sigma_t^2)}$，则频域核为高斯核：
\[
\widehat{w}(\omega) \approx \frac{\sqrt{2\pi}\,\sigma_t}{2T}\, e^{-\sigma_t^2 \omega^2 / 2}.
\]
\end{corollary}

\begin{corollary}[经典极限：矩形窗与 $\delta$ 核]\label{cor:classical_limit}
取矩形窗 $w(t) = \mathbf{1}_{[-T,T]}(t)$，则 $\widehat{w}(\omega) = \frac{\sin(\omega T)}{\omega T}$，且当 $T \to \infty$ 时 $\widehat{w}(\omega_m - \omega_n) \to \delta_{mn}$。
\end{corollary}

\begin{corollary}[零时延与余弦相似度]\label{cor:zero_delay_cosine}
取 $\tau = 0$ 并归一化，则
\[
\frac{(f \star g)(0)}{\|f\| \|g\|}
\]
恰为 $L^2([0,T])$ 中的余弦相似度。
\end{corollary}

\begin{property}[相位卷积、振幅卷积、频率卷积、趋势卷积的基本性质]\label{prop:convolution_types_properties}
\begin{enumerate}[label=(\roman*)]
    \item $|\mathbb{P}_{mn}| = 1$，$\mathbb{P}_{mn} = \overline{\mathbb{P}'_{nm}}$，其中 $\mathbb{P}'_{mn} := e^{i(\varphi_m - \theta_n)}$ 为交换 $f$ 与 $g$ 的角色后所得的相位卷积（需配合相位定义域约定，见注 \ref{rem:phase_convolution_asymmetry}）；
    \item $\mathbb{A}_{mn} = \mathbb{A}_{nm} \ge 0$；
    \item $\mathbb{F}_{mn} = \mathbb{F}_{nm}$；
    \item $\mathbb{T}_{mn} = \mathbb{A}_{mn} \mathbb{P}_{mn} = \overline{\mathbb{T}_{nm}}$；
    \item 时间反转使差角族与和角族互换；
    \item 振幅开关 $\alpha$ 从 $0$ 到 $1$ 即相位卷积到趋势卷积。
\end{enumerate}
\end{property}

\begin{proof}
逐条由定义直接验证。
\end{proof}

\begin{example}[周期为 2 的相位/振幅/频率卷积矩阵]\label{ex:period_two_matrix}
设 $T = 2$，$f(t) = c_0 + c_1 e^{i\pi t}$，$g(t) = d_0 + d_1 e^{i\pi t}$。则：
\[
\mathbb{P} = \begin{pmatrix}
e^{i(\theta_0 - \varphi_0)} & e^{i(\theta_0 - \varphi_1)} \\
e^{i(\theta_1 - \varphi_0)} & e^{i(\theta_1 - \varphi_1)}
\end{pmatrix}, \quad
\mathbb{A} = \begin{pmatrix}
A_0 B_0 & A_0 B_1 \\
A_1 B_0 & A_1 B_1
\end{pmatrix}, \quad
\mathbb{T} = \begin{pmatrix}
c_0 \overline{d_0} & c_0 \overline{d_1} \\
c_1 \overline{d_0} & c_1 \overline{d_1}
\end{pmatrix}.
\]
频率卷积矩阵为 $\mathbb{F} = \begin{pmatrix} 0 & 0 \\ 0 & \pi^2 \end{pmatrix}$（$\omega_0 = 0$，$\omega_1 = \pi$）。
\end{example}

\begin{example}[河图相位卷积的具体矩阵]\label{ex:hetu_phase_convolution}
取 $T = 10$，$f = g$ 为河图复变函数 $f(z) = \sum_{n=0}^{9} c_n e^{2\pi i n z/10}$，其中 $c_n = A_n e^{i\phi_n}$。相位卷积矩阵 $\mathbb{P}$ 为 $10 \times 10$ 矩阵，元素 $\mathbb{P}_{mn} = e^{i(\phi_m - \phi_n)}$。按 $\mathbb{Z}_{10} \cong \mathbb{Z}_2 \times \mathbb{Z}_5$ 分解，$\mathbb{P}$ 分成四个 $5 \times 5$ 块：
\[
\mathbb{P} = \begin{pmatrix} \mathbb{P}^{(\text{阳},\text{阳})} & \mathbb{P}^{(\text{阳},\text{阴})} \\ \mathbb{P}^{(\text{阴},\text{阳})} & \mathbb{P}^{(\text{阴},\text{阴})} \end{pmatrix},
\]
其中 $(\text{阳},\text{阳})$ 块对应 $m, n \in \{0,1,2,3,4\}$，$(\text{阴},\text{阴})$ 块对应 $m, n \in \{5,6,7,8,9\}$，交叉块对应阴阳混合的相位交互。
\end{example}

\subsection{\texorpdfstring{周期卷积：四参数算子的完整交互}{周期卷积}}
\label{subsec:periodic_convolution_series}

\begin{definition}[周期卷积]\label{def:periodic_convolution}
设 $f, g$ 为周期函数（或概周期函数），其周期级数展开为
\[
f(t) = \sum_{p \in P} \sum_{n \in \mathbb{Z}} A_{p,n} e^{i(\omega_{p,n} t + \phi_{p,n})}, \qquad
g(t) = \sum_{q \in Q} \sum_{m \in \mathbb{Z}} B_{q,m} e^{i(\nu_{q,m} t + \psi_{q,m})},
\]
其中 $\omega_{p,n} = 2\pi n / T_p$，$\nu_{q,m} = 2\pi m / T_q$。定义\textbf{频率合成算子}与\textbf{相位合成算子}为
\[
\Omega_{p,q,n,m} := \omega_{p,n} + \nu_{q,m}, \qquad
\Phi_{p,q,n,m} := \phi_{p,n} + \psi_{q,m} \pmod{2\pi}.
\]
$f$ 与 $g$ 的\textbf{周期卷积}定义为
\[
(f \circledast g)(t) = \sum_{p \in P} \sum_{q \in Q} \sum_{n, m \in \mathbb{Z}} A_{p,n} B_{q,m} \cdot \operatorname{Similarity}^{(d)}(\theta_{p,n}, \theta_{q,m}) \cdot e^{i(\Omega_{p,q,n,m} t + \Phi_{p,q,n,m})},
\]
其中：
\begin{itemize}
    \item $\theta_{p,n} = (\omega_{p,n}, \phi_{p,n})$ 为 $f$ 的第 $(p,n)$ 分量的相位点；
    \item $\theta_{q,m} = (\nu_{q,m}, \psi_{q,m})$ 为 $g$ 的第 $(q,m)$ 分量的相位点；
    \item $\operatorname{Similarity}^{(d)}$ 为 $d$ 维相位相似核（$d = 1, \dots, 6$，对应 6 个复相位核，命题 \ref{prop:six_complex_kernels}）。
\end{itemize}
周期卷积以\textbf{振幅}（缩放）、\textbf{频率}（平移）、\textbf{相位}（旋转）、\textbf{周期}（定长）四个参数为算子，保留全部频率对与全部周期对的交互。当 $P = Q = \{T\}$（单周期层）且仅保留对角项 $n = m$ 时，周期卷积退化为经典卷积（定义 \ref{def:convolution}）。
\end{definition}

\begin{proposition}[周期卷积的退化情形]\label{prop:periodic_convolution_degenerations}
\begin{enumerate}[label=(\roman*)]
    \item \textbf{单周期层 + 对角核}：$P = Q = \{T\}$，$W \to \delta_{nm}$，$S^{(1)}$（差角），$\alpha = 1$，则周期卷积退化为经典卷积：
    \[
    (f \circledast g)(t) \to (f * g)(t) = \sum_n c_n d_n e^{i\omega_n t}.
    \]
    \item \textbf{单周期层 + 全矩阵}：$P = Q = \{T\}$，$W \equiv 1$，$S^{(1)}$，$\alpha = 1$，则周期卷积退化为趋势卷积：
    \[
    (f \circledast g)_{n,m} = \mathbb{T}_{nm} = c_n \overline{d_m}.
    \]
    \item \textbf{单周期层 + 单位模}：$P = Q = \{T\}$，$W \equiv 1$，$S^{(1)}$，$\alpha = 0$，则周期卷积退化为相位卷积：
    \[
    (f \circledast g)_{n,m} = \mathbb{P}_{nm} = e^{i(\theta_n - \varphi_m)}.
    \]
\end{enumerate}
\end{proposition}

\begin{proof}
(i) 当 $W \to \delta_{nm}$ 时，非对角项消失，仅保留 $n = m$ 项。由 $S^{(1)}(n,n) = 1$ 和 $\Phi_{n,n} = \theta_n + \varphi_n$，得 $\sum_n A_n B_n e^{i(2\omega_n t + \theta_n + \varphi_n)} = \sum_n c_n d_n e^{i\omega_n t}$。
(ii) $W \equiv 1$ 时全部频率对保留，$\alpha = 1$ 保留振幅乘积，得趋势卷积矩阵 $\mathbb{T}_{nm} = c_n \overline{d_m}$。
(iii) $W \equiv 1$，$\alpha = 0$ 丢弃振幅，得相位卷积矩阵 $\mathbb{P}_{nm} = e^{i(\theta_n - \varphi_m)}$。
\end{proof}

\begin{theorem}[周期卷积的函子性]\label{thm:periodic_convolution_functoriality_series}
周期卷积 $\circledast$ 是平移周期范畴上的双函子，满足：
\begin{enumerate}[label=(\roman*)]
    \item \textbf{保持恒等}：$\operatorname{id}_f \circledast \operatorname{id}_g = \operatorname{id}_{f \circledast g}$。
    \item \textbf{保持复合}：$(h_2 \circ h_1) \circledast (k_2 \circ k_1) = (h_2 \circledast k_2) \circ (h_1 \circledast k_1)$。
    \item \textbf{与相位卷积兼容}：$\mathbb{P}(f \circledast g) \cong \mathbb{P}(f) \otimes \mathbb{P}(g)$。
    \item \textbf{与趋势卷积兼容}：$\mathbb{T}(f \circledast g) \cong \mathbb{T}(f) \otimes \mathbb{T}(g)$。
\end{enumerate}
\end{theorem}

\begin{proof}
(i)-(ii) 由周期卷积的线性性和傅里叶变换的函子性保证。(iii)-(iv) 由周期卷积的频域表示（外积 $\mathbf{c}_f \otimes \mathbf{c}_g$）与相位/趋势卷积的定义直接验证。周期卷积在频域中是张量积，相位与趋势在频域中是逐点运算。
\end{proof}

\begin{proposition}[周期卷积与四类经典卷积的关系]\label{prop:periodic_convolution_hierarchy}
周期卷积是卷积理论五级阶梯的顶点，各级阶梯之间的退化关系为：
\[
\text{周期卷积} \circledast \xrightarrow{\text{单周期层}} \text{趋势卷积} \mathbb{T} \xrightarrow{\text{对角核 } W \to \delta} \text{经典卷积} *,
\]
\[
\text{周期卷积} \circledast \xrightarrow{\text{单周期层} + \alpha = 0} \text{相位卷积} \mathbb{P}.
\]
此外，趋势卷积 $=$ 振幅卷积 $\times$ 相位卷积（逐项），频率卷积 $=$ 频率级数外积，三者从振幅、相位、频率三个独立分量刻画卷积。
\end{proposition}

\begin{proof}
由命题 \ref{prop:periodic_convolution_degenerations} 与 \ref{prop:three_convolutions_relation} 直接得出。
\end{proof}

\begin{remark}[五级阶梯的完整谱系]\label{rem:five_stage_hierarchy}
卷积理论的五级阶梯为：\textbf{经典卷积}（只保留同频率对 $m = n$ 的时延乘积）$\to$ \textbf{相位卷积/振幅卷积/频率卷积}（分别保留全部频率对的相位、振幅、频率分量）$\to$ \textbf{趋势卷积}（保留全部频率对的振幅与相位）$\to$ \textbf{周期级数}（多周期层分解）$\to$ \textbf{周期卷积}（四参数算子，保留全部频率对与全部周期对的交互）。五级阶梯的每一级都是上一级的退化，且每一级的信息量严格递增。
\end{remark}

\begin{figure}[htbp]
\centering
\begin{tikzpicture}[scale=0.9,>=stealth,
  node distance=1.1cm,
  ladder/.style={draw, rounded corners, minimum width=11.0cm, minimum height=0.8cm, align=center, font=\small},
  arrow/.style={->, thick, >=stealth}]

  \node[ladder, fill=blue!10] (s1) {经典卷积 $f * g = \sum_n c_n d_n e^{i\omega_n t}$（仅同频率对）};
  \node[ladder, fill=cyan!12, below=of s1] (s2) {相位卷积 / 振幅卷积 / 频率卷积（全部频率对的单分量投影）};
  \node[ladder, fill=green!12, below=of s2] (s3) {趋势卷积 $\mathbb{T}_{mn} = c_m \overline{d_n}$（全部频率对的振幅与相位）};
  \node[ladder, fill=orange!12, below=of s3] (s4) {周期级数（多周期层分解，四元组参数化）};
  \node[ladder, fill=red!15, below=of s4] (s5) {周期卷积 $f \circledast g$（四参数算子，全部频率对与全部周期对）};

  \draw[arrow] (s1) -- node[right, font=\scriptsize, text=gray!70!black, align=left] {退化} (s2);
  \draw[arrow] (s2) -- node[right, font=\scriptsize, text=gray!70!black, align=left] {综合} (s3);
  \draw[arrow] (s3) -- node[right, font=\scriptsize, text=gray!70!black, align=left] {分层} (s4);
  \draw[arrow] (s4) -- node[right, font=\scriptsize, text=gray!70!black, align=left] {四参数化} (s5);

  \node[font=\scriptsize, gray!70!black, align=center, text width=0.8\textwidth] at (7.5,-3.6) {信息量严格递增};
\end{tikzpicture}
\caption{卷积理论五级阶梯。从经典卷积（最简，仅同频率对）到周期卷积（顶点，四参数化），每一级都是上一级的退化或综合，信息量严格递增。周期卷积同时保留振幅、频率、相位与周期四个分量的全部交互信息。}
\label{fig:convolution_five_stages}
\end{figure}

\begin{property}[周期卷积的核心参数化]\label{prop:periodic_convolution_parametrization}
周期卷积（定义 \ref{def:periodic_convolution}）以四个参数为算子：
\begin{enumerate}[label=(\roman*)]
  \item \textbf{振幅}（缩放）：通过振幅开关 $\alpha$ 控制是否保留振幅乘积 $A_{p,n} B_{q,m}$；
  \item \textbf{频率}（平移）：通过频率合成算子 $\Omega_{p,q,n,m} = \omega_{p,n} + \nu_{q,m}$ 合成新频率；
  \item \textbf{相位}（旋转）：通过相位相似核 $\operatorname{Similarity}^{(d)}(\theta_{p,n}, \theta_{q,m})$ 保留相位关系；
  \item \textbf{周期}（定长）：通过周期类指标集 $P \times Q$ 保留全部周期对的交互。
\end{enumerate}
四参数的完整保留，使周期卷积成为卷积理论中信息量最大的算子。
\end{property}

\begin{proof}
由定义 \ref{def:periodic_convolution} 中周期卷积的显式公式，四个参数分别出现在四个因子上：振幅 $A_{p,n}B_{q,m}$、频率合成 $\Omega_{p,q,n,m}$、相位核 $\operatorname{Similarity}^{(d)}$、周期指标 $P \times Q$。命题 \ref{prop:periodic_convolution_degenerations} 表明，去除任一参数即退化为下一级阶梯的算子。
\end{proof}

\section*{本章七主题对照表}

\begin{center}
\renewcommand{\arraystretch}{1.3}
\begin{tabular}{>{\centering\arraybackslash}p{2.2cm}>{\centering\arraybackslash}p{10.5cm}}
\toprule
主题 & 本章对应 \\
\midrule
平移 & 泰勒级数的局域平移展开；傅里叶级数的全局平移周期；卷积的平移性质（命题 \ref{prop:convolution_properties}）；河图复变函数的半周期平移（命题 \ref{prop:hetu_yin_yang_flip}） \\
缩放 & 振幅级数（定义 \ref{def:amplitude_series}）；振幅卷积（定义 \ref{def:amplitude_convolution}）；趋势卷积的振幅乘积（定义 \ref{def:tendency_convolution}） \\
旋转 & 相位级数（定义 \ref{def:phase_series}）；相位卷积（定义 \ref{def:phase_convolution}）；6 个复相位核（命题 \ref{prop:six_complex_kernels}）；河图相位算子（定义 \ref{def:hetu_phase_operator}） \\
守恒 & 帕塞瓦尔恒等式（定理 \ref{thm:parseval_fourier_classical}）；振幅级数的帕塞瓦尔（命题 \ref{prop:amplitude_parseval}）；玻尔广义帕塞瓦尔恒等式（定理 \ref{thm:bohr_parseval}） \\
对称 & 洛朗级数的正负幂对称；傅里叶基的正交对称（命题 \ref{prop:fourier_orthogonality}）；振幅卷积与频率卷积的对称性（命题 \ref{prop:convolution_types_properties}）；河图相位算子群的直积分解（命题 \ref{prop:hetu_phase_group}） \\
闭合 & 周期级数系数的唯一性（定理 \ref{thm:periodic_series_bohr}）；周期卷积的函子性闭合格（定理 \ref{thm:periodic_convolution_functoriality_series}） \\
周期 & 频率级数（定义 \ref{def:frequency_series}）；周期级数的四元组结构（定义 \ref{def:periodic_series}）；周期卷积的五级阶梯（命题 \ref{prop:periodic_convolution_hierarchy}）；河图五行分量的周期结构（命题 \ref{prop:hetu_five_periods}） \\
\bottomrule
\end{tabular}
\end{center}

\section*{本章思考题}

\begin{enumerate}[leftmargin=3em, label=\textbf{题 \arabic*}]
    \item 证明比值判别法与根值判别法（命题 \ref{prop:ratio_root_tests}）：
    设 $a_k\neq0$，若
    \[
    \lim_{k\to\infty}\left|\frac{a_{k+1}}{a_k}\right|=\rho
    \quad\text{或}\quad
    \lim_{k\to\infty}\sqrt[k]{|a_k|}=\rho,
    \]
    则 $\rho<1$ 时级数绝对收敛，$\rho>1$ 时发散，$\rho=1$ 时判别失效。
    请分别写出两种判别法的完整证明，并各举一个 $\rho=1$ 时失效的例子。

    \item 证明魏尔斯特拉斯 $M$-判别法（定理 \ref{thm:weierstrass_mtest}）：
    设 $\sum_k M_k$ 收敛，且 $|u_k(t)|\le M_k$ 对所有 $t\in I$ 成立，
    则 $\sum_k u_k(t)$ 在 $I$ 上一致且绝对收敛。
    再证明一致收敛的逐项积分（定理 \ref{thm:uniform_termwise_integration}）与逐项求导
    （定理 \ref{thm:uniform_termwise_differentiation}）。

    \item 求几何级数 $\sum_{k=0}^{\infty}r^k$ 的收敛条件与和函数（例 \ref{ex:geometric_series}），
    并利用它推导：
    \begin{enumerate}[label=(\roman*)]
        \item $\displaystyle \sum_{k=0}^{\infty}kr^k=\frac{r}{(1-r)^2}$（$|r|<1$）；
        \item $\displaystyle \sum_{k=0}^{\infty}k^2r^k=\frac{r(1+r)}{(1-r)^3}$（$|r|<1$）。
    \end{enumerate}

    \item 写出 $e^x$、$\sin x$、$\cos x$、$\ln(1+x)$ 与 $(1+x)^\alpha$ 的泰勒级数
    （例 \ref{ex:taylor_examples}），并利用带拉格朗日余项的泰勒公式
    （定理 \ref{thm:taylor_lagrange}）证明：
    \[
    e^x=\sum_{k=0}^{\infty}\frac{x^k}{k!}
    \]
    对一切 $x\in\mathbb{R}$ 成立。
    进一步说明泰勒展开与傅里叶展开的局域-全局对偶（注 \ref{rem:taylor_fourier_duality}）。

    \item 验证洛朗级数的系数唯一性（引理 \ref{lem:laurent_uniqueness}），
    并写出 $e^{1/z}$ 在 $|z|>0$ 上的洛朗展开（例 \ref{ex:laurent_example}）。
    根据主部的形态，判断 $z_0=0$ 是 $e^{1/z}$ 的哪一类奇点
    （命题 \ref{prop:laurent_singularity}）。

    \item 设 $f$ 以 $T$ 为周期且在 $[0,T]$ 上可积，其复傅里叶系数为 $c_n$。
    证明傅里叶基的正交性（命题 \ref{prop:fourier_orthogonality}）：
    \[
    \frac{1}{T}\int_0^T e^{i2\pi(n-m)t/T}\,dt=\delta_{nm}.
    \]
    并利用正交性证明帕塞瓦尔恒等式（定理 \ref{thm:parseval_fourier_classical}）：
    \[
    \frac{1}{T}\int_0^T|f(t)|^2\,dt=\sum_{n\in\mathbb{Z}}|c_n|^2.
    \]

    \item 写出方波（在 $[0,T/2)$ 上取 $1$，在 $[T/2,T)$ 上取 $-1$）的傅里叶级数
    （例 \ref{ex:fourier_square_wave}）。
    说明狄利克雷收敛定理（定理 \ref{thm:dirichlet_convergence}）在跳跃点处给出的收敛值，
    并解释吉布斯现象（例 \ref{ex:gibbs_phenomenon}）为什么不会随项数增加而消失。

    \item 证明费耶核与一致收敛定理（定理 \ref{thm:fejer_uniform}）：
    设 $f$ 连续且以 $T$ 为周期，费耶平均
    \[
    \sigma_Nf(t)=\frac{1}{N+1}\sum_{k=0}^{N}S_kf(t)
    \]
    在 $\mathbb{R}$ 上一致收敛于 $f$。
    请说明费耶核是好核（非负、积分为 $1$、在原点外一致趋于零），
    并解释为什么费耶平均能避免吉布斯现象。

    \item 证明河图傅里叶分解（定理 \ref{thm:hetu_fourier_decomposition}）：
    设 $f$ 满足 $f(z+10)=f(z)$ 且傅里叶展开一致收敛，则
    \[
    f(z)=\sum_{k=0}^{4}\left(c_k+c_{k+5}e^{\pi iz}\right)e^{2\pi ikz/10}.
    \]
    取具体函数
    \[
    f(z)=3+2e^{\pi iz/5}+e^{\pi iz}+5e^{3\pi iz/5}+4e^{2\pi iz},
    \]
    写出其五个 $\mathbb{Z}_5$ 分量，并验证半周期平移 $z\mapsto z+5$ 使奇数分量变号
    （命题 \ref{prop:hetu_yin_yang_flip}）。

    \item 证明相位级数与时间平移的关系（命题 \ref{prop:phase_translation}）：
    若 $\mathcal{P}_f(t)=\sum_n e^{i(nt+\phi_n)}$，
    则时间平移 $f_\tau(t)=f(t-\tau)$ 对应相位级数各项旋转 $-n\tau$。
    再证明有限相位级数的模长守恒（定理 \ref{thm:phase_modulus_conservation}）：
    \[
    \frac{1}{2\pi}\int_0^{2\pi}|S_N(t)|^2dt=2N+1.
    \]

    \item 证明振幅级数的帕塞瓦尔恒等式（命题 \ref{prop:amplitude_parseval}）：
    设 $f\in L^2([0,2\pi])$，其傅里叶系数为 $c_n=A_ne^{i\phi_n}$，则
    \[
    \frac{1}{2\pi}\int_0^{2\pi}|f(t)|^2dt=\sum_{n\in\mathbb{Z}}A_n^2.
    \]
    并说明能量仅由振幅谱决定，与相位谱无关（推论 \ref{cor:amplitude_energy_only}）。

    \item 证明周期卷积的退化情形（命题 \ref{prop:periodic_convolution_degenerations}）：
    \begin{enumerate}[label=(\roman*)]
        \item 单周期层且 $W\to\delta_{nm}$、$S^{(1)}$、$\alpha=1$ 时，
              周期卷积退化为经典卷积；
        \item 单周期层且 $W\equiv1$、$S^{(1)}$、$\alpha=1$ 时，
              周期卷积退化为趋势卷积 $\mathbb{T}_{nm}=c_n\overline{d_m}$；
        \item 单周期层且 $W\equiv1$、$S^{(1)}$、$\alpha=0$ 时，
              周期卷积退化为相位卷积 $\mathbb{P}_{nm}=e^{i(\theta_n-\varphi_m)}$。
    \end{enumerate}

    \item 综合题：说明振幅级数、相位级数与频率级数分别是周期级数在三个分量上的投影
    （命题 \ref{prop:periodic_series_synthesis} 与推论 \ref{cor:periodic_series_projections}）。
    请以周期级数
    \[
    f(t)=3e^{i(\sqrt{2}t+\pi/6)}+2e^{i(t-\pi/3)}
    \]
    为例，分别写出其振幅级数、相位级数与频率级数的显式形式，
    并指出为什么任意单一投影都不足以还原 $f$。
\end{enumerate}

